%% toc: P3 36 36.1,36.3
\chapter{Mostowski Quantifiers and the Uncountability Quantifier}
\chaptermark{Mostowski Quantifiers and \(Q_1\)}
\label{sec:p03-w03-p03-042-mostowski}
\index{generalized quantifier!Mostowski}

\emph{Purpose.}  This chapter gives Mostowski's original account of
generalized quantifiers, as printed in \emph{Fundamenta Mathematicae}~44
(1957), pp.~12--36 \cite{Mostowski1957}, separately from the
isomorphism-invariant family of \cref{def:p02-generalized-quantifier}, and
then connects it with the uncountability quantifier \(Q_1\) of
\cref{def:p03-q-one}.  It also states and proves the effective half of
Keisler's completeness theorem for \(L(Q)\) and derives the host statement
\cref{thm:p03-keisler-q-one} from Keisler's proved theorems.
\emph{Inputs}: a set \(I\) and functions on its subsets for Mostowski's
quantifiers; a first-order language \(L\) with at most countably many
symbols for Keisler's logic.  \emph{Conclusions}: Mostowski's
classification of quantifiers by cardinal pairs, his classification of
the quantifiers on a denumerable set modulo the numerical ones, his
results on the completeness problem, his form of the
Skolem--L\"owenheim theorem for quantifiers, his undecidability and
decidability results for the monadic calculus, the recursive
enumerability of the valid sentences of \(L(Q)\), and the exact form of the host theorem on
\(L(Q_1)\).

Mostowski's paper has six sections: propositional functions and
quantifiers (\S1), examples (\S2), a formal calculus and its
interpretation (\S3), the completeness problem (\S4), the
Skolem--L\"owenheim theorem (\S5), and the monadic calculus (\S6).
All six are presented below, \S6 in
\cref{sec:p03-w03-p03-042-monadic}.  Throughout this part Mostowski's
notation is kept; the bridge to the handbook's notation is
\cref{sec:p03-w03-p03-042-bridge}.

\section{Quantifiers limited to a set}
\label{sec:p03-w03-p03-042-limited}

\begin{hbconvention}{def:p03-w03-p03-042-truth-values}{Mostowski's truth values and cardinal notation}
Mostowski writes \(\vee\) for the truth value ``truth'' and \(\wedge\) for
``falsity'' \cite[p.~12]{Mostowski1957}.  The same signs, together with
\(\sim\), also denote join, meet and complement in the two-element Boolean
algebra \(\{\vee,\wedge\}\) and in every Boolean algebra used below, so that,
for instance, \({\sim}\vee=\wedge\) and \(\vee\wedge\wedge=\wedge\).  Inside
the formulas of the calculus~(S) of
\cref{def:p03-w03-p03-042-calculus-s} the signs \(\vee,\wedge\) are
connectives; the context always fixes which role is meant.  For a set
\(X\), \(\overline{X}\) is its cardinal number, and
\(\underset{x\in X}{E}[W(x)]\) is the set of elements \(x\in X\)
satisfying \(W(x)\); for \(f:X\to Y\), \(f^{-1}(y)\) is the set of
\(x\in X\) with \(f(x)=y\) \cite[p.~13, note~4]{Mostowski1957}.
\end{hbconvention}

\begin{hbdefinition}{def:p03-w03-p03-042-prop-function}{Propositional functions on a set (Mostowski, \S1)}
Let \(I\) be a set and \(I^{*}=I\times I\times\cdots\) the set of infinite
sequences \((x_1,x_2,\ldots)\) with \(x_j\in I\).  A map
\(F:I^{*}\to\{\vee,\wedge\}\) is a \emph{propositional function on \(I\)}
if there is a finite set \(K\) of integers such that \(F(x)=F(y)\)
whenever \(x,y\in I^{*}\) and \(x_j=y_j\) for all \(j\in K\).  The smallest
such \(K\) is the \emph{support} of \(F\).  A propositional function whose
support has exactly one element \(j\) is said to have \emph{one
argument}; it is identified with the subset of those \(a\in I\) for which
\(F(x)=\vee\) whenever \(x_j=a\), and this subset is again written
\(F^{-1}(\vee)\).  For a one-one map \(\varphi\) of \(I\) onto a set
\(I'\), \(\varphi(x)=(\varphi(x_1),\varphi(x_2),\ldots)\), and \(F_\varphi\)
is the propositional function on \(I'\) with
\(F_\varphi(\varphi(x))=F(x)\) \cite[pp.~12--13]{Mostowski1957}.
\end{hbdefinition}

The support exists because the finite sets \(K\) with the stated property
are closed under intersection (see \cref{thm:p03-w03-p03-042-support}).
Under the identification, a one-argument propositional function is a
subset \(X\subseteq I\), and \(F_\varphi\) is the image \(\varphi[X]\).

\begin{hblemma}{thm:p03-w03-p03-042-support}{proved=annex:p03-w03-p03-042-support}{Supports exist}
If \(K_1\) and \(K_2\) are finite sets of integers such that \(F(x)=F(y)\)
whenever \(x\) and \(y\) agree on \(K_1\), and likewise for \(K_2\), then
the same holds for \(K_1\cap K_2\).  Hence every propositional function
has a least such set.
\end{hblemma}

\begin{hbdefinition}{def:p03-w03-p03-042-limited-quantifier}{Quantifiers limited to \(I\) and unlimited quantifiers (Mostowski, \S1)}
A \emph{quantifier limited to \(I\)} is a function \(Q\) assigning one of
\(\vee,\wedge\) to every propositional function \(F\) on \(I\) with one
argument, such that
\[
 Q(F)=Q(F_\varphi)\qquad\text{for every such }F\text{ and every
 permutation }\varphi\text{ of }I
\]
\cite[p.~13]{Mostowski1957}.  An \emph{unlimited quantifier} (or simply a
\emph{quantifier}) is a function \(Q\) assigning to each set \(I\) a
quantifier \(Q_I\) limited to \(I\) such that \(Q_I(F)=Q_{I'}(F_\varphi)\)
for every propositional function \(F\) on \(I\) with one argument and every
one-one map \(\varphi\) of \(I\) onto \(I'\) \cite[p.~13]{Mostowski1957}.
Boolean operations on limited and on unlimited quantifiers are taken
pointwise, for example \((Q'_I\vee Q''_I)(F)=Q'_I(F)\vee Q''_I(F)\)
\cite[p.~14]{Mostowski1957}.
\end{hbdefinition}

Mostowski reads the first clause as generalizing the fact that the
classical quantifiers turn one-argument propositional functions into
propositions, and the invariance clause as the requirement that a
quantifier cannot tell different elements of \(I\) apart
\cite[p.~13]{Mostowski1957}.  The arity is fixed at one: each quantifier
binds a single variable and acts on one subset of \(I\).

\begin{hbdefinition}{def:p03-w03-p03-042-qt}{The quantifier \(Q_T\) and duals (Mostowski, \S1)}
Let \((\mathfrak m_\xi,\mathfrak n_\xi)\) be the sequence of all pairs of
cardinal numbers with \(\mathfrak m_\xi+\mathfrak n_\xi=\overline{I}\).
For a function \(T\) assigning a truth value to each such pair, put
\[
 Q_T(F)=T\bigl(\overline{F^{-1}(\vee)},\,\overline{F^{-1}(\wedge)}\bigr),
\]
where \(F^{-1}(\vee)\) and \(F^{-1}(\wedge)\) are the subsets of \(I\) on
which the one-argument function \(F\) takes the value \(\vee\), resp.\
\(\wedge\) \cite[p.~13]{Mostowski1957}.  If \(Q=Q_T\), \(T\)
\emph{determines} \(Q\).  With \(T^{*}(\mathfrak m,\mathfrak n)=
{\sim}T(\mathfrak n,\mathfrak m)\), the quantifier \(Q_{T^{*}}\) is the
\emph{dual} of \(Q_T\), written \(Q_T^{*}\) \cite[p.~13]{Mostowski1957}.
\end{hbdefinition}

\begin{hbtheorem}{thm:p03-w03-p03-042-theorem-one}{proved=annex:p03-w03-p03-042-theorem-one}{Mostowski's Theorem~1: quantifiers are functions of cardinal pairs}
\(Q_T\) is a quantifier limited to \(I\).  For each quantifier \(Q\)
limited to \(I\) there is a \(T\) with \(Q_T=Q\), and this \(T\) is unique
\cite[p.~13]{Mostowski1957}.  Moreover \(Q_T^{*}(F)={\sim}Q_T({\sim}F)\)
for every one-argument \(F\); this last identity is added here to make the
meaning of the dual explicit.
\end{hbtheorem}

Mostowski calls the theorem easily provable and gives no proof; the annex
supplies one.  Existence of \(F\) with prescribed cardinals uses only the
definition of the cardinal sum: from \(\mathfrak m+\mathfrak n=
\overline I\) one obtains disjoint sets of powers \(\mathfrak m,\mathfrak
n\) whose union is equipotent with \(I\).  Mostowski states no set-theoretic
framework; treating cardinal numbers as objects that can be compared and
listed, as in the sequence of pairs above, is where the axiom of choice
enters his paper.

\begin{hbexample}{def:p03-w03-p03-042-examples}{Mostowski's examples of quantifiers (\S2)}
\begin{enumerate}
\item[(a)] If \(T(\mathfrak m,\mathfrak n)=\vee\) iff \(\mathfrak m\neq0\),
      then \(Q_T\) is the existential quantifier \(\exists\) limited to
      \(I\), and its dual is the general quantifier \(\forall\) limited to
      \(I\) \cite[p.~14]{Mostowski1957}.
\item[(b)] For non-negative integers \(m,n\), let \(T'(\mathfrak m,
      \mathfrak n)=\vee\) iff \(\mathfrak m=m\) and \(T''(\mathfrak m,
      \mathfrak n)=\vee\) iff \(\mathfrak n=n\).  Then \(Q_{T'}\) and
      \(Q_{T''}\) are written \(\overset{(m)}{\underset{I}{\Sigma}}\) and
      \(\overset{(n)}{\underset{I}{\Pi}}\); the unlimited quantifiers
      assigning them to each \(I\) are \(\overset{(m)}{\Sigma}\) and
      \(\overset{(n)}{\Pi}\).  Boolean polynomials in the quantifiers
      \(\overset{(m)}{\Sigma},\overset{(n)}{\Pi}\) (\(m,n=0,1,2,\ldots\))
      are the \emph{numerical quantifiers}.  For instance, with
      \(Q^{(1)}=\overset{(m_1)}{\Sigma}\vee\cdots\vee
      \overset{(m_k)}{\Sigma}\) and \(Q^{(2)}=\overset{(m_1)}{\Pi}\vee\cdots\vee
      \overset{(m_k)}{\Pi}\), on an infinite \(I\) the value
      \(Q^{(1)}_I(F)=\vee\) (resp.\ \(Q^{(2)}_I(F)=\vee\)) says that the set
      of \(x\) with \(F(x,\ldots)=\vee\) (resp.\ \(=\wedge\)) has exactly
      \(m_1\) or \(\ldots\) or exactly \(m_k\) elements
      \cite[p.~14]{Mostowski1957}.
\item[(c)] \(S_I\) is \(Q_{T_1}\) with \(T_1(\mathfrak m,\mathfrak n)=\vee\)
      iff \(\mathfrak m<\aleph_0\), and \(S^0_I\) is \(Q_{T_2}\) with
      \(T_2(\mathfrak m,\mathfrak n)=\vee\) iff \(\mathfrak m<\aleph_0\) or
      \(\mathfrak n<\aleph_0\); \(S\) and \(S^0\) are the corresponding
      unlimited quantifiers.  Thus \(S_I(F)=\vee\) says that at most
      finitely many \(x\) have \(F(x,\ldots)=\vee\), and \(S^0_I(F)=\vee\)
      says that at most finitely many \(x\) have \(F(x,\ldots)=\vee\) or at
      most finitely many have \(F(x,\ldots)=\wedge\)
      \cite[p.~14]{Mostowski1957}.  For the printed form of \(T_2\) see
      \cref{def:p03-w03-p03-042-misprint}.
\item[(d)] \(P_I\) is \(Q_T\) with \(T(\mathfrak m,\mathfrak n)=\vee\) iff
      \(\mathfrak m\le\aleph_0\); \(P\) is the unlimited quantifier.
      \(P_I(F)=\vee\) says that at most denumerably many \(x\in I\) have
      \(F(x,\ldots)=\vee\) \cite[pp.~14--15]{Mostowski1957}.
\end{enumerate}
On a finite \(I\), \(S_I\) and \(S^0_I\) are the constant quantifier with
value \(\vee\); on a denumerable \(I\), so is \(P_I\).
\end{hbexample}

\begin{hbwarning}{def:p03-w03-p03-042-misprint}{A misprint in Mostowski's definition of \(S^0\)}
The printed display on p.~14 defines \(T_2\) by
``\(\{T_2(\mathfrak m_\xi,\mathfrak n_\xi)=\wedge\}\equiv(\mathfrak m_\xi<
\aleph_0\text{ or }\mathfrak n_\xi<\aleph_0)\)'', with \(\wedge\), while the
sentence after it reads \(S^0_I(F)=\vee\) as ``finitely many \(\vee\)-values
or finitely many \(\wedge\)-values'' \cite[p.~14]{Mostowski1957}.  The two
are complementary.  The verbal reading is the one used everywhere else in
the paper: in \S2(e) \(S^0_I\) corresponds to \(\langle N,N,\wedge\rangle\),
which requires \(T_2(n,\aleph_0)=T_2(\aleph_0,n)=\vee\) and
\(T_2(\aleph_0,\aleph_0)=\wedge\) \cite[p.~15]{Mostowski1957}, and the
formula used for \(S^0\) in the proof of Theorem~3 must hold in the ring of
integers, where every set \(\{z:0<z<x\}\) is finite
\cite[p.~19]{Mostowski1957}.  This chapter therefore uses
\(T_2(\mathfrak m,\mathfrak n)=\vee\) iff \(\mathfrak m<\aleph_0\) or
\(\mathfrak n<\aleph_0\), as in item~(c) above.
\end{hbwarning}

\begin{hbtheorem}{thm:p03-w03-p03-042-denumerable-classification}{proved=annex:p03-w03-p03-042-denumerable-classification}{Quantifiers on a denumerable set modulo numerical quantifiers (Mostowski, \S2(e))}
Let \(I\) be denumerable and \(N=\{0,1,2,\ldots\}\).
\begin{enumerate}
\item A quantifier \(Q_T\) limited to \(I\) is determined by the values
      \(T(n,\aleph_0)\), \(T(\aleph_0,n)\) (\(n=0,1,2,\ldots\)) and
      \(\gamma_T=T(\aleph_0,\aleph_0)\).
\item Let \(\mathcal A\) be the set of quantifiers \(Q_T\) limited to \(I\)
      for which \(T(n,\aleph_0)\) and \(T(\aleph_0,n)\) are constant from
      some \(n\) on, \(\mathfrak I\) the ideal of finite subsets of \(N\),
      \(\mathfrak A\) the Boolean algebra of subsets of \(N\) generated by
      \(\mathfrak I\), and \(\mathfrak C=\mathfrak A\times\mathfrak A\times
      \{\wedge,\vee\}\).  Then \(\mathcal A\) is a Boolean algebra and
      \(Q_T\mapsto\langle A_T,B_T,\gamma_T\rangle\), with
      \(A_T=\underset{n}{E}[T(n,\aleph_0)=\vee]\) and
      \(B_T=\underset{n}{E}[T(\aleph_0,n)=\vee]\), is an isomorphism of
      \(\mathcal A\) onto \(\mathfrak C\).  \(S_I\), \(S_I^{*}\), \(S_I^0\)
      correspond to \(X_1=\langle N,0,\wedge\rangle\),
      \(X_2=\langle N,0,\vee\rangle\), \(X_3=\langle N,N,\wedge\rangle\),
      where \(0\) is the void set.
\item \(\mathfrak I_0=\mathfrak I\times\mathfrak I\times\{\wedge\}\) is an
      ideal of \(\mathfrak C\); the quantifiers whose images lie in
      \(\mathfrak I_0\) form an ideal \(\mathcal A_0\) of \(\mathcal A\), and
      all of them are numerical.  Every element of \(\mathfrak C\) is
      congruent modulo \(\mathfrak I_0\) to one of the eight elements
      \(X_0=\langle0,0,\wedge\rangle,X_1,X_2,X_3,{\sim}X_0,{\sim}X_1,
      {\sim}X_2,{\sim}X_3\).
\item Hence every \(Q_T\in\mathcal A\) is either numerical or congruent
      modulo \(\mathcal A_0\) to one of \(S_I,S^{*}_I,S^0_I,{\sim}S_I,
      {\sim}S^{*}_I,{\sim}S^0_I\); and if \(Q_T\in\mathcal A\) is not
      numerical, then one of \(S_I,S^{*}_I,S^0_I\) is a Boolean polynomial
      in \(Q_T\) and numerical quantifiers, so is definable from them
      \cite[p.~15]{Mostowski1957}.
\end{enumerate}
\end{hbtheorem}

Mostowski credits the proof of the last item to J.~\L{}o\'s
\cite[p.~15, note~6]{Mostowski1957}; the paper states the conclusions
with ``it is easy to see'', and the annex writes out each step.

\section{The calculus (S), satisfaction and definability}
\label{sec:p03-w03-p03-042-calculus}

\begin{hbdefinition}{def:p03-w03-p03-042-calculus-s}{The calculus (S) and its interpretation (Mostowski, \S3)}
(S) is the first-order functional calculus with identity, with individual
variables \(\mathrm x_1,\mathrm x_2,\ldots\) and functional (predicate)
variables \(\mathrm F_1,\mathrm F_2,\ldots\) of given degrees, except that
instead of the symbols for the existential and general quantifiers it has
\(s\) symbols \(\mathbf Q^1,\ldots,\mathbf Q^s\): if \(\mathrm F\) is a
formula and \(\mathrm x\) a variable, \((\mathbf Q^j\mathrm x)\mathrm F\)
is a formula (\(j=1,\ldots,s\)), and \(\mathrm x\) is bound in it.  A
formula with only bound individual variables is \emph{closed}
\cite[p.~15]{Mostowski1957}.  The primitive connective is the stroke
\(\mathrm Z_1|\mathrm Z_2\).

Let \(Q^1,\ldots,Q^s\) be unlimited quantifiers and \(I\) a set.  An
\emph{\(I\)-valuation} \(M\) assigns to each individual variable
\(\mathrm x_i\) an element \([\mathrm x_i]_M\in I\) and to each functional
variable \(\mathrm F_j\) of degree \(k\) a propositional function
\([\mathrm F_j]_M\) on \(I\) with support \(\{1,\ldots,k\}\).  The map
\(\mathit{val}_{MI}\) from formulas to \(\{\vee,\wedge\}\) is defined by
induction \cite[p.~16]{Mostowski1957}:
\begin{enumerate}
\item \(\mathit{val}_{MI}(\mathrm F(\mathrm x_{i_1},\ldots,\mathrm
      x_{i_k}))=[\mathrm F]_M([\mathrm x_{i_1}]_M,\ldots,[\mathrm
      x_{i_k}]_M,[\mathrm x_{i_k}]_M,[\mathrm x_{i_k}]_M,\ldots)\);
\item \(\mathit{val}_{MI}(\mathrm x_i=\mathrm x_j)=\vee\) iff
      \([\mathrm x_i]_M=[\mathrm x_j]_M\);
\item \(\mathit{val}_{MI}(\mathrm Z_1|\mathrm Z_2)={\sim}\mathit{val}_{MI}
      (\mathrm Z_1)\vee{\sim}\mathit{val}_{MI}(\mathrm Z_2)\);
\item \(\mathit{val}_{MI}((\mathbf Q^j\mathrm x_i)\mathrm Z_1)=Q^j_I(F)\),
      where \(F\) is the one-argument propositional function
      \(\underset{x\in I}{E}[\mathit{val}_{M(i,x)I}(\mathrm Z_1)=\vee]\) and
      \(M(i,x)\) differs from \(M\) only in assigning \(x\) to
      \(\mathrm x_i\).
\end{enumerate}
\(\mathrm Z\) is \emph{true in \(I\)} (\emph{satisfiable in \(I\)}) if
\(\mathit{val}_{MI}(\mathrm Z)=\vee\) for every (for at least one) \(M\);
\(\mathrm Z\) is \emph{true} (\emph{satisfiable}) if it is true in every
\(I\) (satisfiable in at least one \(I\)) \cite[p.~16]{Mostowski1957}.
These notions are relative to the interpretation of each \(\mathbf Q^j\)
as \(Q^j\), which is held fixed \cite[pp.~16--17]{Mostowski1957}.
\end{hbdefinition}

The quantifier clause is printed with ``support \(\{i\}\)'' and the
argument written \(y_1\) \cite[p.~16]{Mostowski1957}; under the
identification of one-argument functions with subsets of \(I\) the two
readings define the same set, which is the one displayed in clause~4.
The other connectives and \(\exists,\forall\) are not primitive; when
\(\exists\) and \(\forall\) are among \(Q^1,\ldots,Q^s\) they are the
quantifiers of \cref{def:p03-w03-p03-042-examples}(a), and
\(\supset,\wedge,\vee,\equiv\) abbreviate stroke combinations.

\begin{hblemma}{thm:p03-w03-p03-042-val-invariance}{proved=annex:p03-w03-p03-042-val-invariance}{Free variables and isomorphism invariance of \(\mathit{val}\)}
\begin{enumerate}
\item \(\mathit{val}_{MI}(\mathrm Z)\) depends only on the values
      \([\mathrm x_i]_M\), \([\mathrm F_j]_M\) of the variables free in
      \(\mathrm Z\) \cite[p.~16]{Mostowski1957}.
\item If \(\varphi\) maps \(I\) one-one onto \(I'\) and \(M'\) is the
      \(I'\)-valuation with \([\mathrm x_i]_{M'}=\varphi([\mathrm x_i]_M)\)
      and \([\mathrm F_j]_{M'}=([\mathrm F_j]_M)_\varphi\), then
      \(\mathit{val}_{M'I'}(\mathrm Z)=\mathit{val}_{MI}(\mathrm Z)\) for
      every formula \(\mathrm Z\).
\end{enumerate}
\end{hblemma}

Mostowski states part~1 as easy to prove; part~2 is used without
statement in his proofs of Theorems~2 and~3, where a formula holds in all
copies of the ring of integers once it holds in one.

\begin{hbdefinition}{def:p03-w03-p03-042-definability}{Definability of quantifiers (Mostowski, \S3)}
\(Q^1_I\) is \emph{definable in \(I\) in terms of} \(Q^2_I,\ldots,Q^s_I\)
if there is a formula \(\mathrm Z\) of (S) not containing \(\mathbf Q^1\),
with \(\mathrm F\) as its unique free variable, such that
\((\mathbf Q^1\mathrm x)\mathrm F(\mathrm x)\equiv\mathrm Z\) is true in
\(I\).  If this holds for each \(I\) with the same \(\mathrm Z\), the
unlimited quantifier \(Q^1\) is \emph{definable in terms of}
\(Q^2,\ldots,Q^s\) \cite[p.~17]{Mostowski1957}.
\end{hbdefinition}

A Boolean polynomial in limited quantifiers is definable from them in this
sense: \(\mathrm Z\) is the corresponding stroke combination of formulas
\((\mathbf Q^j\mathrm x)\mathrm F(\mathrm x)\).  This is the sense of
``definable'' in the last item of
\cref{thm:p03-w03-p03-042-denumerable-classification}.

\section{The completeness problem}
\label{sec:p03-w03-p03-042-completeness-problem}

\begin{hbdefinition}{def:p03-w03-p03-042-property-e}{The completeness problem and property (E) (Mostowski, \S4)}
The \emph{completeness problem} for \(Q^1,\ldots,Q^s\) (for
\(Q^1_I,\ldots,Q^s_I\)) asks whether the set of true formulas of (S) (the
set of formulas true in \(I\)) is recursively enumerable; it has a
\emph{positive} solution if it is, a \emph{negative} one if it is not
\cite[p.~17]{Mostowski1957}.  Formulas are identified with their numbers
in a fixed effective numbering of the expressions of (S), which
Mostowski leaves implicit.  A quantifier
\(Q\) has \emph{property (E)} if, for each denumerable \(I\), the function
\(T\) determining \(Q_I\) makes both sets
\(\underset{n}{E}[T(n,\aleph_0)=\wedge]\) and
\(\underset{n}{E}[T(n,\aleph_0)=\vee]\) denumerable
\cite[p.~17]{Mostowski1957}.  For a limited quantifier \(Q_I\) on a
denumerable \(I\) the same condition on its \(T\) is called property (E)
of \(Q_I\).
\end{hbdefinition}

\begin{hbdefinition}{def:p03-w03-p03-042-ring-axioms}{The ring of integers inside (S) (Mostowski, \S4)}
Mostowski takes \(\mathrm A_0\) to be the conjunction of the
axioms of the elementary theory of non-densely ordered rings from Tarski's
\emph{Undecidable Theories} \cite[p.~17]{Mostowski1957}; that list is not
reproduced here.  The proofs use \(\mathrm A_0\) only through facts
(a)--(d) below and the fact that the integers satisfy it, so the following
explicit list is fixed: with \(\mathrm F_1(\mathrm x)\) read ``\(\mathrm x\)
is zero'', \(\mathrm F_2(\mathrm x,\mathrm y,\mathrm z)\) ``\(\mathrm x\) is
the sum of \(\mathrm y\) and \(\mathrm z\)'', \(\mathrm F_3(\mathrm x,
\mathrm y,\mathrm z)\) ``\(\mathrm x\) is the product of \(\mathrm y\) and
\(\mathrm z\)'', and \(\mathrm F_4(\mathrm x,\mathrm y)\) ``\(\mathrm x\) is
less than \(\mathrm y\)'' \cite[p.~17]{Mostowski1957}, \(\mathrm A_0\) says
in first-order terms with \(\exists,\forall\) that exactly one element is
zero, that \(\mathrm F_2\) and \(\mathrm F_3\) are graphs of binary
operations, that these make the universe a commutative ring with a unit
element, that \(\mathrm F_4\) is a strict linear order compatible with
addition (\(y<z\) implies \(y+w<z+w\)) and with multiplication by positive
elements, and that the unit is the least positive element.  Write \(\theta,
\oplus,\odot,\prec\) for the zero, the operations and the order determined
by a valuation \(M\) satisfying \(\mathrm A_0\) in \(I\), and \(1\) for the
unit.  Then (a) there is exactly one \(\theta\); (b) \(\oplus,\odot\) are
operations; (c) \(\prec\) orders \(I\); (d) \(I\) is a non-densely ordered
ring with zero \(\theta\): there is no element strictly between \(y\) and
\(y\oplus1\), since such a \(z\) would give \(\theta\prec z\ominus y\prec
1\).  In particular every \(y\) has the immediate predecessor \(y\ominus1\)
and the positive elements \(1,1\oplus1,\ldots\) are distinct, so \(I\) is
infinite.  \(\mathrm P(\mathrm x)\) is \((\exists\mathrm u)[\mathrm
F_1(\mathrm u)\wedge\mathrm F_4(\mathrm u,\mathrm x)]\), and
\(Z_x=\underset{z\in I}{E}[\theta\prec z\prec x]\).
\end{hbdefinition}

\begin{hblemma}{thm:p03-w03-p03-042-ring-lemma}{proved=annex:p03-w03-p03-042-ring-lemma}{Rings satisfying the axioms}
\begin{enumerate}
\item If \(M\) satisfies \(\mathrm A_0\) in \(I\) and \(Z_x\) is finite
      for every positive \(x\), then \(n\mapsto n\cdot1\) is an isomorphism
      of the ordered ring of integers onto \((I,\theta,\oplus,\odot,\prec)\).
\item The set of formulas of (S) built with \(\exists,\forall\) only and
      with no free variables other than \(\mathrm F_1,\ldots,\mathrm F_4\)
      that are true at the standard valuation \(M_0\) of the integers
      \(I_0\) (that of Mostowski's (f)) is not recursively enumerable.
\end{enumerate}
\end{hblemma}

\begin{hbtheorem}{thm:p03-w03-p03-042-theorem-two}{proved=annex:p03-w03-p03-042-theorem-two}{Mostowski's Theorem~2}
If the general and the existential quantifiers and at least one
quantifier \(Q\) with property (E) occur among \(Q^1,\ldots,Q^s\), then
the completeness problem for these quantifiers has a negative solution
\cite[p.~17]{Mostowski1957}.
\end{hbtheorem}

Mostowski's proof axiomatizes the ring of integers up to isomorphism with
the help of \(Q\) and then uses the fact, which he cites as known, that
the set of true sentences about the ring of integers is not recursively
enumerable, indeed not arithmetically definable
\cite[p.~19]{Mostowski1957}.  The annex derives that fact from
Tarski's undefinability theorem (\cref{thm:tarski-undefinability}) by a
computable translation.

\begin{hbtheorem}{thm:p03-w03-p03-042-theorem-three}{proved=annex:p03-w03-p03-042-theorem-three}{Mostowski's Theorem~3}
If \(\exists,\forall,S\) or \(\exists,\forall,S^0\) occur among
\(Q^1,\ldots,Q^s\), then the completeness problem for these quantifiers,
and the completeness problem for these quantifiers limited to a
denumerable set, both have a negative solution
\cite[p.~19]{Mostowski1957}.
\end{hbtheorem}

Property (E) is not necessary for the negative solution: neither \(S\)
nor \(S^0\) has it, since \(T_1(n,\aleph_0)=T_2(n,\aleph_0)=\vee\) for
all \(n\) \cite[p.~19]{Mostowski1957}.

\begin{hbtheorem}{thm:p03-w03-p03-042-theorem-four}{proved=annex:p03-w03-p03-042-theorem-four}{Mostowski's Theorem~4}
Let \(\exists,\forall\) occur among \(Q^1,\ldots,Q^s\) and let \(I\) be
denumerable.  The completeness problem for \(Q^1_I,\ldots,Q^s_I\) has a
positive solution if and only if \(Q^1_I,\ldots,Q^s_I\) are all numerical
\cite[p.~20]{Mostowski1957}.
\end{hbtheorem}

Mostowski calls sufficiency obvious and proves necessity from Theorems~2
and~3 and \S2(e); the annex writes out both directions.  The general case
of the completeness problem and its case for quantifiers limited to
non-denumerable sets are left open in the paper; Mostowski sees no way of
solving it even for \(\exists,\forall,P\) \cite[p.~20]{Mostowski1957}.
The next theorem shows that the method above does not apply there.

\begin{hbtheorem}{thm:p03-w03-p03-042-theorem-five}{proved=annex:p03-w03-p03-042-theorem-five}{Mostowski's Theorem~5}
If \(Q\) is a quantifier definable in terms of \(\exists,\forall,P\), then
\(Q_I\neq S_I\) and \(Q_I\neq S^0_I\) for each denumerable set \(I\)
\cite[p.~20]{Mostowski1957}.
\end{hbtheorem}

\section{The Skolem--L\"owenheim theorem for quantifiers (Mostowski, \S5)}
\label{sec:p03-w03-p03-042-skolem-lowenheim}

\emph{Purpose.}  Mostowski's \S5 asks for which quantifiers the
Skolem--L\"owenheim theorem survives in the calculus (S) of
\cref{def:p03-w03-p03-042-calculus-s}.  \emph{Inputs}: unlimited
quantifiers \(Q^1,\ldots,Q^s\) interpreting \(\mathbf Q^1,\ldots,\mathbf
Q^s\).  \emph{Conclusions}: a single cardinal condition on the determining
functions is sufficient (Theorem~6) and, when \(\exists,\forall\) are
present, necessary (Theorem~7) for every closed formula satisfiable in an
infinite set to be satisfiable in a denumerable one; and the Tarski form
of the theorem forces all the quantifiers to be definable from
\(\exists,\forall\) set by set (Theorem~8).  The set-theoretic background
is Mostowski's: cardinals are compared and well ordered and sets are well
ordered, so the axiom of choice is used throughout, and it is named in the
proofs where it enters.

For an unlimited quantifier \(Q\) and a set \(I\), write \(T_I\) for the
function determining \(Q_I\) (\cref{def:p03-w03-p03-042-qt}); by
\cref{thm:p03-w03-p03-042-transfer}(1) below, \(T_I\) depends only on the
power \(\overline I\).

\begin{hbdefinition}{def:p03-w03-p03-042-infinite-powers}{Quantifiers that do not distinguish infinite powers (Mostowski, \S5)}
An unlimited quantifier \(Q\) \emph{does not distinguish infinite powers}
if for any two infinite sets \(I_1,I_2\) the functions \(T_1,T_2\)
determining \(Q_{I_1}\) and \(Q_{I_2}\) satisfy
\[
 T_1(n,\overline{I_1})=T_2(n,\overline{I_2}),\qquad
 T_1(\overline{I_1},n)=T_2(\overline{I_2},n)\qquad(n=0,1,2,\ldots),
\]
and \(T_1(\mathfrak m_1,\mathfrak n_1)=T_2(\mathfrak m_2,\mathfrak n_2)\)
whenever \(\mathfrak m_i+\mathfrak n_i=\overline{I_i}\) and
\(\mathfrak m_i,\mathfrak n_i\geq\aleph_0\) (\(i=1,2\))
\cite[p.~20]{Mostowski1957}.  Otherwise \(Q\) \emph{distinguishes
infinite powers}.
\end{hbdefinition}

The quantifiers \(\exists\), \(\forall\), \(S\), \(S^{*}\), \(S^0\) and
all \(\overset{(m)}{\Sigma}\), \(\overset{(n)}{\Pi}\) do not distinguish
infinite powers: on an infinite set each of their determining functions
depends only on which of its two arguments is finite and, when finite, on
its value (\cref{def:p03-w03-p03-042-examples}).  The quantifier \(P\)
distinguishes them: \(T_I(\aleph_0,\aleph_0)=\vee\) for a denumerable
\(I\), while \(T_I(\aleph_1,\aleph_1)=\wedge\) when \(\overline I=\aleph_1\).
So does \(Q_1={\sim}P\) of \cref{def:p03-q-one}.

\begin{hblemma}{thm:p03-w03-p03-042-transfer}{proved=annex:p03-w03-p03-042-transfer}{Determining functions across sets}
\begin{enumerate}
\item If \(Q\) is an unlimited quantifier and \(\overline I=\overline{I'}\),
      then \(T_I=T_{I'}\).
\item If \(R^1,\ldots,R^r\) do not distinguish infinite powers, neither
      does any Boolean polynomial \(\Phi(R^1,\ldots,R^r)\), and its
      determining function on \(I\) is
      \(\Phi(T^1_I,\ldots,T^r_I)\) computed pointwise.
\item If \(R,R'\) do not distinguish infinite powers and \(R_I=R'_I\) for
      one denumerable \(I\), then \(R_J=R'_J\) for every infinite \(J\).
\end{enumerate}
\end{hblemma}

Part~3 is the step Mostowski carries out in formulas (11)--(14) of the
proof of Theorem~8 \cite[pp.~26--27]{Mostowski1957}; it is isolated here
because the proof of Theorem~8 uses it twice.

\begin{hbdefinition}{def:p03-w03-p03-042-improper}{Constants and improper formulas (Mostowski, \S5)}
For a set \(I\), (S) is enlarged by individual constants \(\sigma_x\),
one for each \(x\in I\).  Formulas containing no \(\sigma_x\) are
\emph{proper}; the others are \emph{improper}.  \(\mathit{val}_{MI}\) is
extended to improper formulas by the stipulation \([\sigma_x]_M=x\) for
every \(I\)-valuation \(M\) \cite[p.~21]{Mostowski1957}.  In this
section a formula is called closed when it has no free individual
variables (constants are allowed); functional variables may be free, as
in \cref{def:p03-w03-p03-042-calculus-s}.
\end{hbdefinition}

\begin{hblemma}{thm:p03-w03-p03-042-lemma-a}{proved=annex:p03-w03-p03-042-lemma-a}{Mostowski's Lemma~(a) of \S5}
If \(\mathrm A\) is a proper formula with the free variables
\(\mathrm F_1,\ldots,\mathrm F_k,\mathrm x_1,\ldots,\mathrm x_l\), then
\(\mathit{val}_{MI}(\mathrm A)=\vee\) if and only if
\([\mathrm F_1]_M,\ldots,[\mathrm F_k]_M\) satisfy in \(I\) the improper
formula \(\mathrm A'\) obtained from \(\mathrm A\) by substituting
\(\sigma_{[\mathrm x_j]_M}\) for \(\mathrm x_j\) (\(j=1,\ldots,l\))
\cite[p.~21]{Mostowski1957}.  More generally, for any formula
\(\mathrm A\) and any free variable \(\mathrm x_j\) of it,
\(\mathit{val}_{MI}(\mathrm A)=\mathit{val}_{MI}(\mathrm A(\sigma_{[\mathrm
x_j]_M}/\mathrm x_j))\).
\end{hblemma}

Mostowski says the lemma is easily proved by induction; the second
sentence is the form in which his proof of Lemma~(b) applies it.

\begin{hbmethod}{def:p03-w03-p03-042-hull}{The denumerable set \(I_0\) (Mostowski, \S5)}
Let \(I\) be infinite and \(\overline M\) an \(I\)-valuation.  Let
\(\mathfrak U\) be the set of proper formulas beginning with one of
\(\mathbf Q^1,\ldots,\mathbf Q^s\).
\begin{enumerate}
\item Choose a denumerable \(I_1\subseteq I\).
\item Given a denumerable \(I_k\), list as \(\mathrm V_1,\mathrm
      V_2,\ldots\) all closed formulas obtained from formulas
      \(\mathrm X\in\mathfrak U\) by substituting constants \(\sigma_x\)
      with \(x\in I_k\) for the free individual variables of \(\mathrm X\).
      For \(\mathrm V_j=(\mathbf Q^h\mathrm x)\mathrm W\), with
      \(\mathrm W(\sigma_a)\) the result of substituting \(\sigma_a\) for
      \(\mathrm x\), put
      \[
       J_1=\underset{a\in I}{E}\bigl[\mathit{val}_{\overline MI}
       (\mathrm W(\sigma_a))=\vee\bigr],\qquad
       J_2=\underset{a\in I}{E}\bigl[\mathit{val}_{\overline MI}
       (\mathrm W(\sigma_a))=\wedge\bigr],
      \]
      with powers \(\mathfrak m_1,\mathfrak m_2\).  If both are infinite,
      let \(I_{kj}\) be a denumerable subset of \(I\) with infinitely many
      elements in common with \(J_1\) and with \(J_2\); if \(\mathfrak m_1\)
      is finite (so \(\mathfrak m_2\) is infinite), a denumerable subset
      containing \(J_1\) and infinitely many elements of \(J_2\); if
      \(\mathfrak m_2\) is finite, a denumerable subset containing \(J_2\)
      and infinitely many elements of \(J_1\)
      \cite[p.~21]{Mostowski1957}.
\item Put \(I_{k+1}=I_k\cup\bigcup_{j}I_{kj}\) and
      \(I_0=\bigcup_{k\geq1}I_k\).
\end{enumerate}
For an \(I\)-valuation \(M\), the \(I_0\)-valuation \(M_0\) agrees with
\(M\) on \(\mathrm x_i\) when \([\mathrm x_i]_M\in I_0\), assigns to
\(\mathrm x_i\) an arbitrary element of \(I_0\) otherwise, and assigns to
\(\mathrm F_j\) the restriction of \([\mathrm F_j]_M\) to \(I_0\)
\cite[p.~22]{Mostowski1957}.
\end{hbmethod}

Mostowski sets \(I_{k+1}=\bigcup_jI_{kj}\).  The term \(I_k\) is added
here so that the sets \(I_k\) increase: the proof of Lemma~(b) needs every
finite set of elements of \(I_0\) to lie in a single \(I_k\), and without
the added term that is not guaranteed.  Each list \(\mathrm V_1,\mathrm
V_2,\ldots\) is countable because (S) has countably many proper formulas
and \(I_k\) is denumerable; choosing the sets \(I_{kj}\) and the
enumerations uses the axiom of choice, and \(I_0\) is a countable union of
denumerable sets, hence denumerable.

\begin{hblemma}{thm:p03-w03-p03-042-lemma-b}{proved=annex:p03-w03-p03-042-lemma-b}{Mostowski's Lemma~(b) of \S5}
Assume that none of \(Q^1,\ldots,Q^s\) distinguishes infinite powers, and
let \(I_0\) be built from \(I\) and \(\overline M\) as in
\cref{def:p03-w03-p03-042-hull}.  Let \(M\) be an \(I\)-valuation with
\([\mathrm x_i]_M\in I_0\) for all \(i\) and \([\mathrm F_j]_M=[\mathrm
F_j]_{\overline M}\) for all \(j\).  If \(\mathrm V\) is a proper formula
and \(\mathrm V'\) results from \(\mathrm V\) by substituting symbols
\(\sigma_a\) (\(a\in I_0\)) for some or all of its free individual
variables, then \(\mathit{val}_{MI}(\mathrm V')=\mathit{val}_{M_0I_0}
(\mathrm V')\) \cite[p.~22]{Mostowski1957}.
\end{hblemma}

\begin{hbtheorem}{thm:p03-w03-p03-042-theorem-six}{proved=annex:p03-w03-p03-042-theorem-six}{Mostowski's Theorem~6}
If none of the quantifiers \(Q^1,\ldots,Q^s\) distinguishes infinite
powers, then each closed formula satisfiable in an infinite set is
satisfiable in a denumerable set \cite[p.~20]{Mostowski1957}.
\end{hbtheorem}

The theorem does not require \(\exists,\forall\) among the quantifiers.
Its proof follows Mostowski's: the set \(I_0\) is closed under witnesses
for every quantified subformula with parameters in \(I_0\), in the sense
that each such subformula sees in \(I_0\) a cardinal pair to which its
quantifier assigns the same value as to the pair it sees in \(I\).

\begin{hbtheorem}{thm:p03-w03-p03-042-theorem-seven}{proved=annex:p03-w03-p03-042-theorem-seven}{Mostowski's Theorem~7}
If \(\exists\) and \(\forall\) occur among \(Q^1,\ldots,Q^s\) and at least
one of these quantifiers distinguishes infinite powers, then there are
closed formulas satisfiable in non-denumerable sets but not satisfiable in
denumerable sets \cite[p.~24]{Mostowski1957}.
\end{hbtheorem}

The proof treats three cases, according to which clause of
\cref{def:p03-w03-p03-042-infinite-powers} fails: (a) \(T_1(n,
\overline{I_1})\neq T_2(n,\overline{I_2})\) for some \(n\); (b)
\(T_1(\overline{I_1},n)\neq T_2(\overline{I_2},n)\) for some \(n\);
(c) \(T_1(\mathfrak m_1,\mathfrak n_1)\neq T_2(\mathfrak m_2,\mathfrak
n_2)\) for infinite \(\mathfrak m_i,\mathfrak n_i\) with \(\mathfrak
m_i+\mathfrak n_i=\overline{I_i}\) \cite[p.~24]{Mostowski1957}.  In each
case a formula says that a relation \(\mathrm G\) is a strict linear order
with an element \(\mathrm x_0\) whose set of predecessors has a prescribed
size or contains a limit point, and that \(\mathbf Q\) gives the value
\(\vee\) to that set of predecessors.  A well ordering of a suitable
non-denumerable set satisfies it; in a denumerable set it would force a
value of \(T\) at a denumerable set that minimality of the counterexample
excludes.  Mostowski writes out case~(a) and says case~(b) is similar
\cite[p.~25]{Mostowski1957}; the annex writes out all three.

\begin{hbwarning}{def:p03-w03-p03-042-limit-point}{A limit point below \(\mathrm x_0\) needs a predecessor}
In the formula for case~(c), the conjunct
\((\exists\mathrm y_1)[\mathrm G(\mathrm y_1,\mathrm x_0)\wedge(\forall
\mathrm s)\{\mathrm G(\mathrm s,\mathrm y_1)\supset(\exists\mathrm t)
[\mathrm G(\mathrm s,\mathrm t)\wedge\mathrm G(\mathrm t,\mathrm y_1)]\}]\)
is meant to say that the set of predecessors of \(\mathrm x_0\) has a
limit point and is therefore infinite \cite[p.~25]{Mostowski1957}.  As
printed it is also satisfied when \(\mathrm y_1\) is the least element,
since the universal clause is then vacuous; a finite set of predecessors
then satisfies it.  The annex adds the conjunct \((\exists\mathrm r)
\mathrm G(\mathrm r,\mathrm y_1)\).  The corresponding conjunct for
\(\mathrm y_2\) needs no change, because \(\mathrm G(\mathrm x_0,\mathrm
y_2)\) already gives \(\mathrm y_2\) a predecessor, and the same holds for
the limit point after \(\mathrm x_0\) in case~(a).  For the satisfying
well ordering in case~(c), Mostowski asks only for a segment of power
\(\mathfrak m_2\) with remainder of power \(\mathfrak n_2\); the annex
chooses each part of order type at least \(\omega+1\) so that both parts
contain a limit point with a predecessor.
\end{hbwarning}

Mostowski records Tarski's generalization of the Skolem--L\"owenheim
theorem, which he cites from Skolem: a formula of the classical functional
calculus satisfiable in an infinite set is satisfiable in every infinite
set \cite[p.~25]{Mostowski1957}.  In the handbook this is the combination
of \cref{thm:downward-ls} and \cref{thm:upward-ls}.  It does not extend to
arbitrary quantifiers:

\begin{hbtheorem}{thm:p03-w03-p03-042-theorem-eight}{proved=annex:p03-w03-p03-042-theorem-eight}{Mostowski's Theorem~8}
If \(\exists\) and \(\forall\) occur among \(Q^1,\ldots,Q^s\) and each
formula of (S) satisfiable in an infinite set is satisfiable in every
infinite set, then for each set \(J\) the quantifiers
\(Q^1_J,\ldots,Q^s_J\) are all definable in terms of \(\exists\) and
\(\forall\) \cite[p.~26]{Mostowski1957}.
\end{hbtheorem}

Definability here is definability in \(J\)
(\cref{def:p03-w03-p03-042-definability}).  The proof gives more, as
Mostowski remarks: each \(Q^j_J\) with \(J\) infinite is a Boolean
polynomial in the quantifiers \(\overset{(n)}{\underset J\Sigma}\),
\(\overset{(m)}{\underset J\Pi}\), and the polynomial does not depend on
\(J\) \cite[p.~27]{Mostowski1957}.

\begin{hbexample}{def:p03-w03-p03-042-theorem-eight-boundary}{Theorem~8 does not give definability of the unlimited quantifiers}
Mostowski observes that under the hypotheses of Theorem~8 the unlimited
quantifiers need not be definable in terms of \(\exists,\forall\)
\cite[p.~27]{Mostowski1957}.  His example is the unlimited quantifier
\(Q\) with \(Q_I=\underset I\Sigma\) for infinite \(I\) and
\(Q_I=\underset I\Pi\) for finite \(I\), with \(Q^1=\exists\),
\(Q^2=\forall\) and \(Q^3=\cdots=Q^s=Q\).  The symbols
\(\underset I\Sigma\) and \(\underset I\Pi\) appear there without the
superscripts of \cref{def:p03-w03-p03-042-examples}(b); they are read here
as the existential and the general quantifier limited to \(I\), in the
classical notation \(\Sigma,\Pi\) for these quantifiers.  With that
reading the hypotheses of Theorem~8 hold and \(Q\) is not definable in
terms of \(\exists,\forall\)
(\cref{thm:p03-w03-p03-042-theorem-eight-boundary}).
\end{hbexample}

\begin{hbtheorem}{thm:p03-w03-p03-042-theorem-eight-boundary}{proved=annex:p03-w03-p03-042-theorem-eight-boundary}{The boundary example}
Let \(Q_I=\exists_I\) for infinite \(I\) and \(Q_I=\forall_I\) for finite
\(I\).  Then \(Q\) is an unlimited quantifier; every formula of (S) with
the quantifiers \(\exists,\forall,Q\) that is satisfiable in an infinite
set is satisfiable in every infinite set; and \(Q\) is not definable in
terms of \(\exists\) and \(\forall\).
\end{hbtheorem}

\emph{Limitation.}  Theorems~6--8 concern satisfiability in sets of a
given power, not the completeness problem; Mostowski's open questions of
\S4 (the general case and quantifiers limited to non-denumerable sets,
\cref{thm:p03-w03-p03-042-theorem-four} and the remark after it) are not
touched by them.  Theorem~7 shows that \(Q_1\) of
\cref{def:p03-q-one}, which distinguishes infinite powers, defeats the
downward theorem; the positive results for \(L(Q_1)\) are of a different
kind (\cref{sec:p03-w03-p03-042-keisler-effective}).


\section{The monadic calculus (Mostowski, \S6)}
\label{sec:p03-w03-p03-042-monadic}

\emph{Purpose.}  Mostowski's \S6 restricts the calculus (S) of
\cref{def:p03-w03-p03-042-calculus-s} to one-place functional variables
and asks when the true formulas of the restricted calculus form a
recursive set.  \emph{Inputs}: unlimited quantifiers \(Q^1,\ldots,Q^s\)
interpreting \(\mathbf Q^1,\ldots,\mathbf Q^s\), and for Theorem~10 a
strictly increasing list of cardinals.  \emph{Conclusions}: some
quantifiers make satisfiability in a denumerable set non-recursive even
for monadic formulas (Theorem~9); cardinality quantifiers ``at least
\(\mathfrak m_j\)'' with \(\mathfrak m_1\) equal to \(1\) or infinite give
a recursive set of true monadic formulas, by an explicit decision method
(Theorem~10); and for every choice of quantifiers the true formulas with a
fixed finite number of functional variables form a recursive set, by an
argument that does not produce the decision method (Theorem~11).  The
axiom of choice is used where cardinals are compared or added, as in the
earlier sections.

\begin{hbdefinition}{def:p03-w03-p03-042-monadic-calculus}{The subsystems \((\mathrm S^M)\) and \((\mathrm S^M_n)\) (Mostowski, \S6)}
\((\mathrm S^M)\) is the part of (S) whose formulas contain only monadic
functional variables, that is, functional variables of one argument; for
\(n\geq1\), \((\mathrm S^M_n)\) is the part of \((\mathrm S^M)\) whose
formulas contain no functional variables other than \(\mathrm
F_1,\ldots,\mathrm F_n\) \cite[p.~28]{Mostowski1957}.  In this chapter the
formulas of \((\mathrm S^M)\) and \((\mathrm S^M_n)\) contain no identity
sign; the reason is \cref{def:p03-w03-p03-042-monadic-identity}.  Truth
and satisfiability are those of \cref{def:p03-w03-p03-042-calculus-s}.
\end{hbdefinition}

The atomic formulas of \((\mathrm S^M)\) are therefore the
\(\mathrm F_h(\mathrm x_l)\), and every formula contains at least one
functional variable.  Since an \(I\)-valuation gives each individual
variable a value in \(I\), there are no \(I\)-valuations when \(I\) is
empty; every formula is true in the empty set and none is satisfiable in
it, so only nonempty sets matter for satisfiability and for the failure of
truth.

\begin{hbwarning}{def:p03-w03-p03-042-monadic-identity}{Identity and the monadic calculus}
(S) is introduced as the first-order calculus \emph{with identity}, with
the quantifier symbols changed \cite[p.~15]{Mostowski1957}, and \S6 does
not say that identity is removed \cite[p.~28]{Mostowski1957}.  The
arguments of \S6 nevertheless use only atoms built from functional
variables: the base step of Lemma~(c) in the proof of Theorem~10, called
evident in the source \cite[p.~30]{Mostowski1957}, fails for
\(\mathrm x_1=\mathrm x_2\), and Theorem~11 is false when identity is
present (\cref{thm:p03-w03-p03-042-monadic-identity}).  The identity-free
reading of \cref{def:p03-w03-p03-042-monadic-calculus} is therefore the
one under which Theorems~10 and~11 and their lemmas are stated and proved
here.  Theorem~9 and its corollary hold in both readings, because the
formulas used in its proof contain no identity.  Whether the conclusion of
Theorem~10 extends to monadic formulas with identity is not settled by
Mostowski's argument, and this chapter makes no claim about it.
\end{hbwarning}

\begin{hblemma}{thm:p03-w03-p03-042-monadic-identity}{proved=annex:p03-w03-p03-042-monadic-identity}{What identity would change}
\begin{enumerate}
\item Let \(Q^1,\ldots,Q^s\) be any quantifiers and \(n\geq1\).  No formula
      \(\mathrm U\) built by the stroke from atoms
      \(\mathrm F_h(\mathrm x_1)\), \(\mathrm F_h(\mathrm x_2)\)
      (\(h\leq n\)) and closed formulas makes \(\mathrm x_1=\mathrm
      x_2\equiv\mathrm U\) true in a set with at least two elements.
\item There is a quantifier \(Q\) such that, with \(s=1\) and
      \(Q^1=Q\), neither the set of true formulas of \((\mathrm S^M_1)\)
      with identity allowed nor the set of those true in a denumerable set
      is recursive.
\end{enumerate}
\end{hblemma}

\subsection{Theorem~9: a non-recursive satisfiability problem}

\begin{hbdefinition}{def:p03-w03-p03-042-hset}{The sets \(H(A)\), constituents and the formulas \(\mathrm W_n\) (Mostowski, \S6)}
For a set \(A\) of positive integers, \(H(A)\) is the set of positive
integers \(p\) for which there are \(n_1,\ldots,n_p\in A\), repetitions
allowed, with \(n_1+n_2+\cdots+n_p\in A\) \cite[p.~28]{Mostowski1957}.
For a monadic functional variable \(\mathrm F_j\), \(\mathrm
F_j^0(\mathrm x)\) is \(\mathrm F_j(\mathrm x)\) and \(\mathrm
F_j^1(\mathrm x)\) is \({\sim}\mathrm F_j(\mathrm x)\).  The \(2^n\)
formulas \(\mathrm F_1^{i_1}(\mathrm x)\wedge\cdots\wedge\mathrm
F_n^{i_n}(\mathrm x)\), \((i_1,\ldots,i_n)\in\{0,1\}^n\), are the
\emph{constituents of order \(n\)}; they are ordered as their index
tuples are ordered lexicographically, and in the proof of Theorem~9 the
first \(n\) of them are written \(\mathrm S_1(\mathrm x),\ldots,\mathrm
S_n(\mathrm x)\).  For a quantifier symbol \(\mathbf Q^3\),
\[
 \mathrm W_n=[(\mathbf Q^3\mathrm x)\mathrm S_1(\mathrm x)]\wedge\cdots
 \wedge[(\mathbf Q^3\mathrm x)\mathrm S_n(\mathrm x)]\wedge
 [(\mathbf Q^3\mathrm x)(\mathrm S_1(\mathrm x)\vee\cdots\vee
 \mathrm S_n(\mathrm x))]
\]
\cite[p.~28]{Mostowski1957}.
\end{hbdefinition}

For each \(a\in I\) and each \(I\)-valuation, \(a\) satisfies exactly one
constituent of order \(n\), the one whose index has \(i_h=0\) exactly when
\(a\) belongs to \([\mathrm F_h]_M\).  The constituents of order \(n\)
therefore cut \(I\) into \(2^n\) disjoint \emph{cells}, and every
partition of \(I\) into \(2^n\) labelled pieces arises in this way from a
suitable choice of \([\mathrm F_1]_M,\ldots,[\mathrm F_n]_M\).

\begin{hblemma}{thm:p03-w03-p03-042-mycielski}{proved=annex:p03-w03-p03-042-mycielski}{Mycielski's lemma}
Let \(A_0\) be the set of integers
\[
 1,\ 2,\ 2^3,\ 2^9,\ \ldots,\ 2^{3^n},\ \ldots
\]
If \(A_1,A_2\subseteq A_0\), \(A_1\neq A_2\), \(1\in A_1\) and \(1\in
A_2\), then \(H(A_1)\neq H(A_2)\) \cite[p.~28]{Mostowski1957}.
Consequently there is a set \(A\) of positive integers for which \(H(A)\)
is not recursive.
\end{hblemma}

Mostowski attributes the proof of the lemma to J.~Mycielski and records as
open whether some \emph{recursive} \(A\) has a non-recursive \(H(A)\)
\cite[p.~28, note~10]{Mostowski1957}.  The ``consequently'' is his
counting step: there are only denumerably many recursive sets, and the
lemma gives continuum many distinct sets \(H(A)\).

\begin{hbtheorem}{thm:p03-w03-p03-042-theorem-nine}{proved=annex:p03-w03-p03-042-theorem-nine}{Mostowski's Theorem~9}
There are quantifiers \(Q^1,\ldots,Q^s\) such that the set of those
formulas of \((\mathrm S^M)\) which are satisfiable in a denumerable set
is not recursive \cite[p.~28]{Mostowski1957}.  One choice is
\(s\geq3\), \(Q^1=\forall\), \(Q^2=\exists\) and
\(Q^3=\cdots=Q^s=Q\), where \(Q_I(X)=\vee\) iff \(\overline X\in A\) and
\(A\) is a set as in \cref{thm:p03-w03-p03-042-mycielski}; for this
choice \(\mathrm W_n\) is satisfiable in a denumerable set iff \(n\in
H(A)\) \cite[pp.~28--29]{Mostowski1957}.
\end{hbtheorem}

\begin{hbcorollary}{thm:p03-w03-p03-042-theorem-nine-true}{proved=annex:p03-w03-p03-042-theorem-nine-true}{True formulas in a denumerable set}
For the quantifiers of \cref{thm:p03-w03-p03-042-theorem-nine}, the set
of formulas of \((\mathrm S^M)\) true in a denumerable set is not
recursive \cite[p.~29]{Mostowski1957}.
\end{hbcorollary}

Mostowski leaves open which quantifiers give a recursive set of true
formulas of \((\mathrm S^M)\), and offers Theorem~10 as a family of
nontrivial examples \cite[p.~29]{Mostowski1957}.

\subsection{Theorem~10: cardinality quantifiers}

\begin{hbtheorem}{thm:p03-w03-p03-042-theorem-ten}{proved=annex:p03-w03-p03-042-theorem-ten}{Mostowski's Theorem~10}
Let \(\mathfrak m_1<\mathfrak m_2<\cdots<\mathfrak m_s\) be \(s\)
cardinals with \(\mathfrak m_2,\ldots,\mathfrak m_s\) infinite and
\(\mathfrak m_1\) either \(1\) or infinite.  Let \(Q^1,\ldots,Q^s\) be
quantifiers such that for each \(I\)
\[
 \{Q^j_I(F)=\vee\}\equiv(\overline F\geq\mathfrak m_j),
 \qquad j=1,2,\ldots,s,
\]
where \(\overline F\) is the power of the set of elements at which \(F\)
takes the value \(\vee\) (the source writes it with a double bar).  Then the set of true formulas of
\((\mathrm S^M)\) is recursive \cite[p.~29]{Mostowski1957}.  The decision
procedure (\cref{thm:p03-w03-p03-042-ten-lemma-d}) is uniform: it uses of
the cardinals only their number \(s\) and whether \(\mathfrak m_1=1\).
\end{hbtheorem}

Mostowski presents the theorem as a generalization of L\"owenheim's
classical theorem and refers to Hilbert--Ackermann for his proof
\cite[p.~29, note~11]{Mostowski1957}.  He also notes: if
\(\mathfrak m_1=1\) then \(Q^1=\exists\); if \(\mathfrak m_2=\aleph_0\)
then \(Q^2={\sim}S\); if \(\mathfrak m_3=\aleph_1\) then \(Q^3={\sim}P\)
\cite[p.~29]{Mostowski1957}, with \(S,P\) as in
\cref{def:p03-w03-p03-042-examples}.  The proof runs through four
lemmas.

\begin{hblemma}{thm:p03-w03-p03-042-ten-lemma-a}{proved=annex:p03-w03-p03-042-ten-lemma-a}{Lemma~(a) of Theorem~10: additivity}
The quantifiers of \cref{thm:p03-w03-p03-042-theorem-ten} satisfy, for
every \(I\) and all one-argument \(F,G\) on \(I\),
\[
 Q^j_I(F\vee G)=Q^j_I(F)\vee Q^j_I(G)\quad(\text{additivity}),\qquad
 Q^j_I(\wedge)=\wedge,
\]
where \(\wedge\) as an argument is the function with constant value
\(\wedge\) \cite[p.~29]{Mostowski1957}.
\end{hblemma}

Additivity fails for ``at least \(\mathfrak m\)'' with \(\mathfrak m\)
finite and greater than \(1\): two singletons have a union with two
elements.  This is the role of the hypothesis on \(\mathfrak m_1\).

\begin{hblemma}{thm:p03-w03-p03-042-ten-lemma-b}{proved=annex:p03-w03-p03-042-ten-lemma-b}{Lemma~(b) of Theorem~10}
Let \(Q^j\) satisfy the equations of
\cref{thm:p03-w03-p03-042-ten-lemma-a}.  If \(\mathrm Z_1,\ldots,\mathrm
Z_n\) are formulas of (S) containing the free variable \(\mathrm x_i\)
and \(\mathrm W_1,\ldots,\mathrm W_n\) are formulas of (S) not containing
\(\mathrm x_i\), then
\[
 (\mathbf Q^j\mathrm x_i)[(\mathrm Z_1\wedge\mathrm W_1)\vee\cdots\vee
 (\mathrm Z_n\wedge\mathrm W_n)]\equiv
 [\mathrm W_1\wedge(\mathbf Q^j\mathrm x_i)\mathrm Z_1]\vee\cdots\vee
 [\mathrm W_n\wedge(\mathbf Q^j\mathrm x_i)\mathrm Z_n]
\]
is true \cite[p.~29]{Mostowski1957}.
\end{hblemma}

Mostowski computes both sides at an arbitrary valuation and reduces the
equivalence to \(w_h\wedge Q^j_I(F_h)=Q^j_I(F_h\wedge w_h)\) for the
constant value \(w_h\) of \(\mathrm W_h\), which is immediate when
\(w_h=\vee\) and follows from \(Q^j_I(\wedge)=\wedge\) otherwise
\cite[pp.~29--30]{Mostowski1957}.  Here \(\mathrm W_1,\ldots,\mathrm W_n\)
are the formulas of the lemma, not the formulas \(\mathrm W_n\) of
\cref{def:p03-w03-p03-042-hset}.

\begin{hbdefinition}{def:p03-w03-p03-042-ten-class}{The formulas \(\mathrm A^j_{i_1\ldots i_n}\) and the class \(\mathfrak R(\mathrm x_1,\ldots,\mathrm x_m)\) (Mostowski, \S6)}
For monadic \(\mathrm F_1,\ldots,\mathrm F_n\), put
\[
 \mathrm A^j_{i_1i_2\ldots i_n}=(\mathbf Q^j\mathrm x)[\mathrm
 F_1^{i_1}(\mathrm x)\wedge\cdots\wedge\mathrm F_n^{i_n}(\mathrm x)],
\]
the quantifier \(\mathbf Q^j\) applied to a constituent of order \(n\),
with a fixed bound variable \(\mathrm x\).  \(\mathfrak R(\mathrm
x_1,\ldots,\mathrm x_m)\) is the least class of formulas containing the
formulas \(\mathrm F_h(\mathrm x_l)\) and \(\mathrm A^j_{i_1\ldots i_n}\)
(\(h=1,\ldots,n\); \(j=1,\ldots,s\); \(l=1,\ldots,m\); each
\(i_t\in\{0,1\}\)) and closed under the stroke.  The case \(m=0\) is
allowed, and the class is then written \(\mathfrak R\); its members are
closed formulas built from the \(\mathrm A^j_{i_1\ldots i_n}\) by the
stroke \cite[p.~30]{Mostowski1957}.
\end{hbdefinition}

\begin{hblemma}{thm:p03-w03-p03-042-ten-lemma-c}{proved=annex:p03-w03-p03-042-ten-lemma-c}{Lemma~(c) of Theorem~10: reduction to \(\mathfrak R\)}
Let the quantifiers satisfy the equations of
\cref{thm:p03-w03-p03-042-ten-lemma-a}.  For each formula \(\mathrm Z\)
of \((\mathrm S^M)\) whose free individual variables are among \(\mathrm
x_1,\ldots,\mathrm x_m\) and whose functional variables are among
\(\mathrm F_1,\ldots,\mathrm F_n\), there is a formula \(\mathrm U\) in
\(\mathfrak R(\mathrm x_1,\ldots,\mathrm x_m)\) such that \(\mathrm
Z\equiv\mathrm U\) is true; \(\mathrm U\) is obtained from \(\mathrm Z\)
by an explicit procedure \cite[p.~30]{Mostowski1957}.
\end{hblemma}

Mostowski states the lemma for formulas whose free variables are exactly
\(\mathrm x_1,\ldots,\mathrm x_m\) and whose functional variables are
exactly \(\mathrm F_1,\ldots,\mathrm F_n\); the form above, with
``among'', is what the induction needs and contains his.  His proof is by
induction on the number of quantifier symbols, the essential case being
\((\mathbf Q^j\mathrm x)\mathrm Z_1\) with \(\mathrm Z_1\) in
\(\mathfrak R(\mathrm x_1,\ldots,\mathrm x_m,\mathrm x)\), which he
writes as a disjunction of constituents in \(\mathrm x\) conjoined with
members of \(\mathfrak R(\mathrm x_1,\ldots,\mathrm x_m)\) and then
treats by Lemma~(b) \cite[p.~30]{Mostowski1957}.  The annex follows this
route.

\begin{hbdefinition}{def:p03-w03-p03-042-allowable}{Allowable valuations and property (T) (Mostowski, \S6)}
Let \(\Phi\) be a Boolean polynomial in the variables \(a^j_{i_1\ldots
i_n}\) (\(j=1,\ldots,s\); each \(i_t\in\{0,1\}\)).  A function
\(\gamma\) giving each of these variables one of the values \(\vee,\wedge\)
is an \emph{allowable valuation} if
\[
 \gamma(a^{j+1}_{i_1\ldots i_n})\leq\gamma(a^j_{i_1\ldots i_n})
 \qquad(j=1,\ldots,s-1),
\]
with \(\wedge<\vee\).  There are finitely many allowable valuations.
\(\Phi\) has \emph{property (T)} if every allowable valuation gives it the
value \(\vee\) \cite[p.~31]{Mostowski1957}.

For \(\mathrm U\) in \(\mathfrak R(\mathrm x_1,\ldots,\mathrm x_m)\) let
\(\Phi_{\mathrm U}\) be the Boolean polynomial in the variables
\(a^j_{i_1\ldots i_n}\) and \(b_{h,l}\) (\(h\leq n\), \(l\leq m\)) from
which \(\mathrm U\) results by substituting \(\mathrm A^j_{i_1\ldots
i_n}\) for \(a^j_{i_1\ldots i_n}\) and \(\mathrm F_h(\mathrm x_l)\) for
\(b_{h,l}\).  For a valuation \(\beta\) of the \(b_{h,l}\), the
\emph{cell} \(c_l(\beta)\) is the tuple \((i_1,\ldots,i_n)\) with
\(i_h=0\) iff \(\beta(b_{h,l})=\vee\).  A pair \((\gamma,\beta)\) is
\emph{admissible} if \(\gamma\) is allowable and, in case
\(\mathfrak m_1=1\), \(\gamma(a^1_{c_l(\beta)})=\vee\) for every \(l\leq
m\) and \(\gamma(a^1_{i_1\ldots i_n})=\vee\) for at least one index.
\end{hbdefinition}

Allowability transcribes the fact that a set with at least
\(\mathfrak m_{j+1}\) elements has at least \(\mathfrak m_j\) elements.
Admissibility, which is not in the source, adds the two facts that hold
for \(\mathfrak m_1=1\) in every nonempty set: the cell of an element is
nonempty, and some cell is nonempty.

\begin{hbwarning}{def:p03-w03-p03-042-lemma-d-correction}{Mostowski's Lemma~(d) needs admissibility when \(\mathfrak m_1=1\)}
Mostowski's Lemma~(d) reads: a formula \(\mathrm Z\) of \(\mathfrak R\)
is true iff it results from a Boolean polynomial with property (T) by
substituting \(\mathrm A^j_{i_1\ldots i_n}\) for \(a^j_{i_1\ldots i_n}\)
\cite[p.~31]{Mostowski1957}.  When \(\mathfrak m_1=1\) this fails.  Take
\(s=1\), \(\mathfrak m_1=1\), \(n=1\) and \(\mathrm Z=\mathrm A^1_0\vee
\mathrm A^1_1\), that is, \((\exists\mathrm x)\mathrm F_1(\mathrm x)\vee
(\exists\mathrm x){\sim}\mathrm F_1(\mathrm x)\): it is true in every
nonempty set, hence true, while its polynomial \(a^1_0\vee a^1_1\) takes
the value \(\wedge\) at the allowable valuation with constant value
\(\wedge\).  In the source's proof of the converse, a cell for which no
\(\gamma(a^j_{i_1\ldots i_n})\) is \(\vee\) has no largest such \(j\);
giving it no elements can leave \(I\) empty, and then there is no
\(I\)-valuation \cite[p.~31]{Mostowski1957}.  When \(\mathfrak m_1\) is
infinite such a cell can receive one element and the source's statement
holds for closed formulas.
\cref{thm:p03-w03-p03-042-ten-lemma-d} states the corrected criterion,
extended to formulas with free variables, which Theorem~10 also covers.
\end{hbwarning}

\begin{hblemma}{thm:p03-w03-p03-042-ten-lemma-d}{proved=annex:p03-w03-p03-042-ten-lemma-d}{Lemma~(d) of Theorem~10, corrected}
Under the hypotheses of \cref{thm:p03-w03-p03-042-theorem-ten}, a formula
\(\mathrm U\) of \(\mathfrak R(\mathrm x_1,\ldots,\mathrm x_m)\) is true
iff \(\Phi_{\mathrm U}\) takes the value \(\vee\) at every admissible pair
\((\gamma,\beta)\).  In particular, if \(\mathfrak m_1\) is infinite, a
formula of \(\mathfrak R\) is true iff its polynomial has property (T),
as Mostowski states.
\end{hblemma}

Mostowski concludes that Theorem~10 follows at once from Lemmas~(c)
and~(d), since the set of polynomials with property (T) is recursive
\cite[p.~32]{Mostowski1957}; the annex does the same with the corrected
criterion.

\subsection{Theorem~11: a fixed number of functional variables}

\begin{hbtheorem}{thm:p03-w03-p03-042-theorem-eleven}{proved=annex:p03-w03-p03-042-theorem-eleven}{Mostowski's Theorem~11}
Let \(Q^1,\ldots,Q^s\) be any quantifiers.  For each \(n\geq1\) the set of
true formulas of \((\mathrm S^M_n)\) and the set of formulas of
\((\mathrm S^M_n)\) which are true in any given set \(I\) are recursive
\cite[p.~32]{Mostowski1957}.
\end{hbtheorem}

Theorem~9 shows that the bound \(n\) on the functional variables cannot be
dropped: its formulas \(\mathrm W_n\) use \(n\) of them.  The proof uses
the following notions, with \(\mathrm S_1(\mathrm x),\ldots,\mathrm
S_m(\mathrm x)\) now standing for \emph{all} \(m=2^n\) constituents of
order \(n\) \cite[p.~32]{Mostowski1957}.

\begin{hbdefinition}{def:p03-w03-p03-042-canonical}{Canonical representations and the formulas \(\mathrm A^j_{\mathrm W}\), \(\mathrm A^j_0\) (Mostowski, \S6)}
A \emph{Boolean polynomial} \(\mathrm W(\mathrm x)\) in \(\mathrm
F_1(\mathrm x),\ldots,\mathrm F_n(\mathrm x)\) is a formula built from
these atoms by the stroke.  If it does not vanish identically it has a
\emph{canonical representation} \(\mathrm S_{k_1}(\mathrm x)\vee\cdots
\vee\mathrm S_{k_h}(\mathrm x)\), unique up to the order of the
disjuncts; here it is written with \(k_1<\cdots<k_h\).  Put
\[
 \mathrm A^j_{\mathrm W}=(\mathbf Q^j\mathrm x)[\mathrm S_{k_1}(\mathrm
 x)\vee\cdots\vee\mathrm S_{k_h}(\mathrm x)],\qquad
 \mathrm A^j_0=(\mathbf Q^j\mathrm x)[\mathrm F_1(\mathrm x)\wedge
 {\sim}\mathrm F_1(\mathrm x)]
\]
\cite[p.~32]{Mostowski1957}.  \(\mathrm A^j_{\mathrm W}\) depends only on
the set \(\{k_1,\ldots,k_h\}\), so there are \(s(2^{2^n}-1)\) formulas
\(\mathrm A^j_{\mathrm W}\) and \(s\) formulas \(\mathrm A^j_0\)
\cite[p.~33]{Mostowski1957}.  When \(\mathrm W\) vanishes identically,
\(\mathrm A^j_{\mathrm W}\) means \(\mathrm A^j_0\)
\cite[p.~33]{Mostowski1957}.  For formulas \(\mathrm C\),
\(\mathrm C^0\) is \(\mathrm C\) and \(\mathrm C^1\) is \({\sim}\mathrm
C\).
\end{hbdefinition}

\begin{hblemma}{thm:p03-w03-p03-042-eleven-lemma-a}{proved=annex:p03-w03-p03-042-eleven-lemma-a}{Lemma~(a) of Theorem~11}
If \(\mathrm A\) and \(\mathrm B\) are formulas not containing the free
variable \(\mathrm x\) and \(\mathrm M(\mathrm x)\), \(\mathrm
N(\mathrm x)\) are formulas in which \(\mathrm x\) is free, then for
\(j=1,\ldots,s\)
\[
 \begin{aligned}
 &(\mathbf Q^j\mathrm x)[(\mathrm A\wedge\mathrm M(\mathrm x))\vee
 (\mathrm B\wedge\mathrm N(\mathrm x))]\\
 &\quad\equiv\{({\sim}\mathrm A\wedge{\sim}\mathrm B\wedge\mathrm A^j_0)
 \vee({\sim}\mathrm A\wedge\mathrm B\wedge(\mathbf Q^j\mathrm x)\mathrm
 N(\mathrm x))\\
 &\qquad\vee(\mathrm A\wedge{\sim}\mathrm B\wedge(\mathbf Q^j\mathrm x)
 \mathrm M(\mathrm x))\\
 &\qquad\vee(\mathrm A\wedge\mathrm B\wedge(\mathbf Q^j\mathrm x)
 [\mathrm M(\mathrm x)\vee\mathrm N(\mathrm x)])\}
 \end{aligned}
\]
is true \cite[p.~32]{Mostowski1957}.
\end{hblemma}

\begin{hblemma}{thm:p03-w03-p03-042-eleven-lemma-b}{proved=annex:p03-w03-p03-042-eleven-lemma-b}{Lemma~(b) of Theorem~11}
Let \(\mathrm Z\) be \(\mathrm M_0(\mathrm x)\vee[\mathrm C_1\wedge
\mathrm N_1(\mathrm x)]\vee\cdots\vee[\mathrm C_p\wedge\mathrm
N_p(\mathrm x)]\), where \(\mathrm C_1,\ldots,\mathrm C_p\) do not contain
\(\mathrm x\).  For \((i_1,\ldots,i_p)\in\{0,1\}^p\) let
\(\underset{i_1,\ldots,i_p}{\Sigma}(\mathrm x)\) be \(\mathrm
M_0(\mathrm x)\vee\mathrm N_{k_1}(\mathrm x)\vee\cdots\vee\mathrm
N_{k_t}(\mathrm x)\), where \(k_1,\ldots,k_t\) are all \(k\leq p\) with
\(i_k=0\).  Then
\[
 (\mathbf Q^j\mathrm x)\mathrm Z\equiv\bigvee_{i_1,\ldots,i_p}
 \bigl[\mathrm C_1^{i_1}\wedge\cdots\wedge\mathrm C_p^{i_p}\wedge
 (\mathbf Q^j\mathrm x)\underset{i_1,\ldots,i_p}{\Sigma}(\mathrm x)\bigr]
\]
is true, the disjunction running over the \(2^p\) index tuples
\cite[p.~32]{Mostowski1957}.
\end{hblemma}

The source prints the disjuncts of \(\Sigma\) as \(\mathrm
N_{i_{k_1}},\ldots\); the indices meant are \(k_1,\ldots,k_t\), as the
defining clause shows \cite[p.~32]{Mostowski1957}.  Mostowski proves
Lemma~(b) by induction on \(p\) from Lemma~(a), taking for \(\mathrm A\) a
formula of the form \(\mathrm D\vee{\sim}\mathrm D\) when \(p=1\) and
splitting off \(\mathrm C_p\) in the inductive step
\cite[pp.~32--33]{Mostowski1957}; the annex proves (b) directly from the
satisfaction clause and checks the inductive equivalence he uses.  Neither
lemma needs any property of the quantifiers.

\begin{hblemma}{thm:p03-w03-p03-042-eleven-lemma-c}{proved=annex:p03-w03-p03-042-eleven-lemma-c}{Lemma~(\(\mathrm c_n\)) of Theorem~11}
Let \(\mathfrak R(\mathrm x_1,\ldots,\mathrm x_m)\) now be the least class
containing the formulas \(\mathrm F_h(\mathrm x_l)\) (\(h\leq n\),
\(l\leq m\)), the \(s\) formulas \(\mathrm A^j_0\) and the
\(s(2^{2^n}-1)\) formulas \(\mathrm A^j_{\mathrm W}\), and closed under the
stroke \cite[p.~33]{Mostowski1957}.  For each formula \(\mathrm Z\) of
\((\mathrm S^M_n)\) whose free individual variables are among \(\mathrm
x_1,\ldots,\mathrm x_m\) there is, by an explicit procedure, a formula
\(\mathrm U\) in \(\mathfrak R(\mathrm x_1,\ldots,\mathrm x_m)\) with
\(\mathrm Z\equiv\mathrm U\) true.
\end{hblemma}

Mostowski reduces (\(\mathrm c_n\)) to the formulas \((\mathbf Q^j\mathrm
x)\mathrm Z_1\) with \(\mathrm Z_1\) either \(\mathrm F_1(\mathrm x)\wedge
{\sim}\mathrm F_1(\mathrm x)\), where \(\mathrm U=\mathrm A^j_0\), or a
disjunction \([\mathrm C_1\wedge\mathrm S_{k_1}(\mathrm x)]\vee\cdots\vee
[\mathrm C_p\wedge\mathrm S_{k_p}(\mathrm x)]\) with the \(\mathrm C_h\)
free of \(\mathrm x\), where Lemma~(b) applies
\cite[p.~33]{Mostowski1957}.

\begin{hbdefinition}{def:p03-w03-p03-042-allowable-for-i}{Valuations allowable for a set \(I\) (Mostowski, \S6)}
Let \(T^j\) be the function determining \(Q^j_I\)
(\cref{def:p03-w03-p03-042-qt}).  A function \(\gamma\) from the
\(s\cdot2^{2^n}\) formulas \(\mathrm A^j_0\), \(\mathrm A^j_{\mathrm W}\)
to truth values is \emph{allowable for \(I\)} if some assignment of
cardinals \(\mathfrak m_k\) to the constituents \(\mathrm S_k(\mathrm x)\)
(\(k=1,\ldots,m\)) satisfies
\begin{enumerate}
\item[\(1^\circ\)] \(\mathfrak m_1+\mathfrak m_2+\cdots+\mathfrak
      m_m=\overline I\);
\item[\(2^\circ\)] if \(\mathrm W(\mathrm x)\) is a Boolean polynomial in
      \(\mathrm F_1(\mathrm x),\ldots,\mathrm F_n(\mathrm x)\) and
      \begin{equation*}
       \mathrm S_{k_1}(\mathrm x)\vee\cdots\vee\mathrm S_{k_p}(\mathrm x),
       \qquad
       \mathrm S_{l_1}(\mathrm x)\vee\cdots\vee\mathrm S_{l_q}(\mathrm x)
       \tag{15}
      \end{equation*}
      are the canonical representations of \(\mathrm W(\mathrm x)\) and
      \({\sim}\mathrm W(\mathrm x)\), then \(\gamma(\mathrm A^j_{\mathrm
      W})=T^j\bigl(\sum_{i=1}^p\mathfrak m_{k_i},\sum_{i=1}^q\mathfrak
      m_{l_i}\bigr)\)
\end{enumerate}
\cite[p.~33]{Mostowski1957}.  A sum in (15) is absent, and the
corresponding cardinal is \(0\), when \(\mathrm W\) or \({\sim}\mathrm
W\) vanishes identically \cite[p.~33]{Mostowski1957}.  For the extension
to free variables, \(\mathfrak K_I\) is the set of pairs \((\gamma,D)\)
with \(D\subseteq\{1,\ldots,m\}\) for which one assignment satisfies
\(1^\circ\), \(2^\circ\) and \(\mathfrak m_k\neq0\) for every \(k\in D\).
\end{hbdefinition}

The cardinals \(\mathfrak m_k\) here are local to Theorem~11 and are not
the cardinals of Theorem~10.  For \(D=\varnothing\) the pairs
\((\gamma,\varnothing)\in\mathfrak K_I\) are exactly the valuations
allowable for \(I\); the set \(\mathfrak K_I\) is finite, as a set of
pairs of functions on finite sets.

\begin{hblemma}{thm:p03-w03-p03-042-eleven-lemma-d}{proved=annex:p03-w03-p03-042-eleven-lemma-d}{Lemma~(d) of Theorem~11}
If \(I\) is nonempty and \(M\) is an \(I\)-valuation, then
\(\gamma(\mathrm A^j_{\mathrm W})=\mathit{val}_{MI}(\mathrm A^j_{\mathrm
W})\) is an allowable valuation for \(I\) \cite[p.~33]{Mostowski1957}.
Moreover, for every \(r\geq0\), \((\gamma,D_r)\in\mathfrak K_I\), where
\(D_r\) is the set of those \(k\) for which \([\mathrm x_l]_M\) satisfies
\(\mathrm S_k\) for some \(l\leq r\).
\end{hblemma}

\begin{hblemma}{thm:p03-w03-p03-042-eleven-lemma-e}{proved=annex:p03-w03-p03-042-eleven-lemma-e}{Lemma~(e) of Theorem~11}
For each valuation \(\gamma\) allowable for a nonempty \(I\) there is an
\(I\)-valuation \(M\) such that \(\gamma(\mathrm A^j_{\mathrm W})=
\mathit{val}_{MI}(\mathrm A^j_{\mathrm W})\) for each \(\mathrm W\)
\cite[p.~34]{Mostowski1957}.  Moreover, if \((\gamma,D)\in\mathfrak K_I\)
and \(c\) maps \(\{1,\ldots,r\}\) into \(D\), \(M\) can be chosen with
\([\mathrm x_l]_M\) satisfying \(\mathrm S_{c(l)}\) for \(l\leq r\).
\end{hblemma}

The parts after ``moreover'' are added here so that formulas with free
individual variables are covered; Mostowski's proof of the theorem
treats closed formulas \cite[p.~35]{Mostowski1957}.  In the source, the
proof of Lemma~(d) uses the sets \(R_k\) of elements satisfying
\(\mathrm S_k\) and their powers, and the proof of Lemma~(e) partitions
\(I\) into sets of the prescribed powers and reads off \([\mathrm
F_i]_M\) as the union of the pieces whose constituents lie in \(\mathrm
F_i\); this is display~(16) \cite[p.~34]{Mostowski1957}.

\emph{Non-effectiveness.}  Mostowski's proof of Theorem~11 shows that the
set of formulas of \(\mathfrak R\) true in \(I\) is the set of those taking
the value \(\vee\) under every valuation in a certain finite set, and
the set of true formulas is obtained in the same way from the union
\(\mathfrak K\) of the finite sets \(\mathfrak K_I\)
\cite[p.~35]{Mostowski1957}.  Which finite set it is depends on the
quantifiers and is not computed by the proof.  Mostowski stresses that
the theorem therefore has no practical use, and illustrates this with the
following example \cite[p.~35]{Mostowski1957}.

\begin{hbtheorem}{thm:p03-w03-p03-042-fermat}{proved=annex:p03-w03-p03-042-fermat}{Mostowski's Fermat-prime example}
Let \(Q^1,Q^2\) be quantifiers such that for each infinite \(I\) and
\(X\subseteq I\), \(Q^1_I(X)=\vee\) iff \(\overline X\) is a
prime and \(Q^2_I(X)=\vee\) iff \(\overline X\) has the form
\(2^{2^n}+1\) \cite[p.~35]{Mostowski1957}.  With the quantifiers
\(\overset{(m)}{\Sigma}\) of \cref{def:p03-w03-p03-042-examples}(b)
\cite[p.~35, note~12]{Mostowski1957}, the formula
\[
 \begin{aligned}
 &[(\mathbf Q^1\mathrm x)\mathrm F(\mathrm x)\wedge(\mathbf Q^2\mathrm x)
 \mathrm F(\mathrm x)]\supset[(\overset{(3)}{\Sigma}\mathrm x)\mathrm
 F(\mathrm x)\vee(\overset{(5)}{\Sigma}\mathrm x)\mathrm F(\mathrm x)\\
 &\qquad\vee(\overset{(17)}{\Sigma}\mathrm x)\mathrm F(\mathrm x)\vee
 (\overset{(257)}{\Sigma}\mathrm x)\mathrm F(\mathrm x)\vee
 (\overset{(65537)}{\Sigma}\mathrm x)\mathrm F(\mathrm x)]
 \end{aligned}
\]
of \((\mathrm S^M_1)\), with \(\mathrm F=\mathrm F_1\) and \(s=7\), is true
in an infinite set \(I\) iff every prime
of the form \(2^{2^n}+1\) is one of \(3,5,17,257,65537\).
\end{hbtheorem}

A prime is a natural number, so an infinite power is never a prime.  The
truth of the formula in an infinite set is thus the question whether
there are more than five Fermat primes, which Mostowski calls famous
\cite[p.~35]{Mostowski1957}.  By Theorem~11 the set of formulas of
\((\mathrm S^M_1)\) true in \(I\) is recursive, yet this does not settle
the question: deciding the formula would need the recursion equations of
the characteristic function, and the proof shows only that such equations
exist \cite[p.~35]{Mostowski1957}.  By contrast, the decision method of
Theorem~10 is given explicitly by Lemmas~(c) and~(d).

\emph{Limitation.}  All three theorems concern the identity-free monadic
calculus (\cref{def:p03-w03-p03-042-monadic-identity}).  Theorem~9 gives
no recursive quantifier with a non-recursive problem: it leaves open
whether \(A\) can be taken recursive.  Theorem~11 decides nothing in
practice, and Theorem~10 covers only the quantifiers ``at least
\(\mathfrak m\)'' with \(\mathfrak m\) equal to \(1\) or infinite.


\section{Quantifiers from partially ordered prefixes (Mundici's overview, \S1)}
\label{sec:p03-w03-p03-042-mundici-po}

\emph{Purpose.}  Mundici's chapter ``Other quantifiers: an overview'' in
the volume \emph{Model-Theoretic Logics} surveys quantifiers introduced
after Mostowski as Lindstr\"om-style classes of structures
\cite[p.~211]{Mundici1985OtherQuantifiers}; this section follows the 1985
Springer printing.  Its first section turns nonlinear prefixes of
\(\forall\) and \(\exists\) into quantifiers.  \emph{Inputs}: first-order
formulas and patterns of pairs of positive integers.
\emph{Conclusions}: the partially ordered quantifiers \(Q_\pi\), the
logic \(\mathscr L^{\circledR}\) in which every projective class of
first-order logic is elementary, and the definability of ``infinitely
many'' with the Henkin quantifier.  Mundici writes \(\mathscr L(Q^i)_{i\in
I}\) for \(\mathscr L_{\omega\omega}(Q^i)_{i\in I}\) and
\(\mathscr L^{\mathrm{mII}}\) for second-order logic with quantifiers over
unary relations only \cite[p.~211]{Mundici1985OtherQuantifiers}.  A class
\(K\) of \(\tau\)-structures is \emph{EC} in a logic \(\mathscr L\) if
\(K=\mathrm{Mod}_{\mathscr L}\psi\) for a sentence \(\psi\) of
\(\mathscr L\), and \emph{PC} in \(\mathscr L_{\omega\omega}\) if it is the
class of \(\tau\)-reducts of the models of a first-order sentence of a
vocabulary \(\tau'\supseteq\tau\).

In a prenex first-order formula an existentially quantified variable
depends on all universally quantified variables before it.  Mundici's
first example \cite[p.~212]{Mundici1985OtherQuantifiers} is the two-row
prefix
\[
 \begin{Bmatrix}\forall x,x'\ \exists y\\ \forall t\ \ \exists z,z'
 \end{Bmatrix}\psi(x,x',t,y,z,z'),
\]
read as \(\exists f,g,g'\,\forall x,x',t\;\psi(x,x',t,f(x,x'),g(t),
g'(t))\), and the second is Henkin's prefix
\[
 \begin{Bmatrix}\forall x\ \exists y\\ \forall x'\ \exists y'
 \end{Bmatrix}\psi(x,x',y,y'),\quad\text{that is,}\quad
 \exists f,f'\,\forall x,x'\;\psi(x,x',f(x),f'(x')).
\]

\begin{hbdefinition}{def:p03-w03-p03-042-mundici-po-quantifier}{Partially ordered quantifiers (Mundici, Definition~1.1.1)}
Let \(\pi=\{\langle n_1,m_1\rangle,\ldots,\langle n_r,m_r\rangle\}\) be a
sequence of pairs of natural numbers \(\geq1\) and
\(s=n_1+m_1+\cdots+n_r+m_r\).  The \emph{partially ordered quantifier with
pattern \(\pi\)} is
\[
 Q_\pi=\{\langle A,R\rangle\mid R\subseteq A^s\text{ and }R\supseteq
 f_1\times\cdots\times f_r\text{ for some }f_1:A^{n_1}\to A^{m_1},\ldots,
 f_r:A^{n_r}\to A^{m_r}\},
\]
and \(Q_\pi\) has \(r\) \emph{rows} \cite[p.~212]{Mundici1985OtherQuantifiers}.
Here a function \(f_i\) is identified with its graph, a subset of
\(A^{n_i+m_i}\), and \(f_1\times\cdots\times f_r\subseteq A^s\) is the
product of the graphs.  The \emph{Henkin quantifier} is
\(Q^{\mathrm H}=Q_\pi\) for \(\pi=\{\langle1,1\rangle,\langle1,1\rangle\}\):
\(Q^{\mathrm H}=\{\langle A,R\rangle\mid R\subseteq A^4\text{ and }R\supseteq
f\times g\text{ for some }f,g:A\to A\}\); the quantifier of the first
example has pattern \(\{\langle2,1\rangle,\langle1,2\rangle\}\)
\cite[p.~212]{Mundici1985OtherQuantifiers}.
\end{hbdefinition}

\begin{hbdefinition}{def:p03-w03-p03-042-mundici-po-logic}{The logics \(\mathscr L(Q_\pi)\) and \(\mathscr L^{\circledR}\) (Mundici, \S1.1)}
In \(\mathscr L(Q)\) with \(Q=Q_\pi\), the quantifier binds \(s\) distinct
variables and is displayed as
\[
 Q=\begin{Bmatrix}\forall x^1_1\cdots x^1_{n_1}\ \exists y^1_1\cdots
 y^1_{m_1}\\ \cdots\qquad\cdots\qquad\cdots\\ \forall x^r_1\cdots
 x^r_{n_r}\ \exists y^r_1\cdots y^r_{m_r}\end{Bmatrix};
\]
the formation rules generalize those for \(\mathscr L(Q^{\mathrm H})\), and
the semantics is read off the display: the existentially quantified
variables in a row depend only on the universally quantified variables of
the same row \cite[p.~213]{Mundici1985OtherQuantifiers}.  Explicitly,
\(\mathfrak A\models Q\chi\) under an assignment iff there are
\(f_i:A^{n_i}\to A^{m_i}\) (\(i=1,\ldots,r\)) such that \(\chi\) holds
whenever each \(\bar y^i\) takes the value \(f_i(\bar x^i)\), the other
free variables keeping their values; this is membership in \(Q_\pi\) of
the relation that \(\chi\) defines in the variables \(\bar x^1\bar
y^1\cdots\bar x^r\bar y^r\).  \(\mathscr L^{\circledR}\) is the smallest
logic in which all partially ordered prefixes of this form are allowed
\cite[p.~213]{Mundici1985OtherQuantifiers}.
\end{hbdefinition}

Mundici reads these quantifiers as expressing uniformization: \(\forall x
\exists y\,Rxy\) says that the binary relation \(R\) can be uniformized,
and the Henkin prefix applied to a 4-ary \(S\) says that \(S\) contains the
product of two uniformizable binary relations
\cite[pp.~212--213]{Mundici1985OtherQuantifiers}.  The Henkin quantifier is
also the subject of \cref{sec:p03-w03-p03-044-henkin-quantifier}, where it
is presented after Ebbinghaus and Zdanowski.

\begin{hbtheorem}{thm:p03-w03-p03-042-mundici-pc-ec}{proved=annex:p03-w03-p03-042-mundici-pc-ec}{Projective classes are elementary in \(\mathscr L^{\circledR}\) (Mundici, Theorem~1.1.2)}
For an arbitrary class \(K\), if \(K\) is PC in \(\mathscr L_{\omega\omega}\)
then \(K\) is EC in \(\mathscr L^{\circledR}\).  Indeed,
\(K=\mathrm{Mod}_{\mathscr L^{\circledR}}\psi\) for some \(\psi\) of the
form \(Q\chi\), where \(Q\) is a partially ordered quantifier as displayed
in \cref{def:p03-w03-p03-042-mundici-po-logic} and \(\chi\in\mathscr
L_{\omega\omega}\) is quantifier-free
\cite[p.~213]{Mundici1985OtherQuantifiers}.
\end{hbtheorem}

Mundici's proof replaces relations by characteristic functions,
Skolemizes, and then normalizes the terms of the matrix
\cite[p.~213]{Mundici1985OtherQuantifiers}.  The annex follows that route
and adds two steps the printed proof leaves implicit: one-element
structures, where characteristic functions cannot separate two truth
values, and Skolem constants, which a row with \(n_i\geq1\) cannot bind
directly.  Skolemization uses the axiom of choice.

The theorem bounds what \(\mathscr L^{\circledR}\) can fail to do: any
property expressible by existentially quantifying relations over a
first-order sentence is a single \(\mathscr L^{\circledR}\)-sentence.  In
particular \(\mathscr L(Q^{\mathrm H})\) is at least as strong as
\(\mathscr L(Q_0)\), where \(Q_0\) is ``there exist infinitely many'':

\begin{hbtheorem}{thm:p03-w03-p03-042-mundici-q-zero}{proved=annex:p03-w03-p03-042-mundici-q-zero}{\(Q_0\) in \(\mathscr L(Q^{\mathrm H})\) (Mundici, \S1.3)}
For every formula \(\varphi(x)\) (with parameters),
\begin{multline*}
Q_0x\,\varphi(x)\quad\text{iff}\quad\exists t\{\varphi(t)\wedge\exists f,g\,
\forall u,v[(u=v\leftrightarrow f(u)=g(v))\\
\wedge(\varphi(u)\to\varphi(f(u))\wedge f(u)\neq t)]\}
\end{multline*}
\cite[p.~216]{Mundici1985OtherQuantifiers}.  The right side is the Henkin
prefix \(\{\forall u\,\exists y;\ \forall v\,\exists y'\}\) applied to
\((u=v\leftrightarrow y=y')\wedge(\varphi(u)\to\varphi(y)\wedge y\neq t)\),
so \(\mathscr L(Q^{\mathrm H})\geq\mathscr L(Q_0)\).  The axiom of choice
is used for the direction from left to right.
\end{hbtheorem}

\subsection{Comparing logics with the Henkin quantifier (Mundici, \S\S1.2--1.4)}
\label{sec:p03-w03-p03-042-mundici-compare}

Mundici's comparisons use the following notions, which he takes from the
general framework of the same volume.  The readings fixed here are the
ones the proofs below use.

\begin{hbconvention}{def:p03-w03-p03-042-mundici-comparison}{Comparison, compactness and axiomatizability of logics}
A \emph{logic} \(\mathscr L\) assigns to each vocabulary \(\tau\) a class
of sentences \(\mathscr L(\tau)\) and a satisfaction relation between
\(\tau\)-structures and these sentences that is invariant under
isomorphism; \(\mathrm{Mod}_{\mathscr L}\psi\) is the class of models of
\(\psi\).
\begin{enumerate}
\item \(\mathscr L\leq\mathscr L'\) means that for every vocabulary
      \(\tau\) and every \(\varphi\in\mathscr L(\tau)\) there is
      \(\psi\in\mathscr L'(\tau)\) with \(\mathrm{Mod}_{\mathscr
      L}\varphi=\mathrm{Mod}_{\mathscr L'}\psi\); \(\mathscr
      L\equiv\mathscr L'\) means \(\mathscr L\leq\mathscr L'\) and
      \(\mathscr L'\leq\mathscr L\).  For the logics \(\mathscr L(Q)\) of
      this chapter the translations constructed below act on formulas and
      keep free variables, which gives the sentence form at once.
\item \(\mathscr L\) is \((\omega,\omega)\)-\emph{compact} if every
      countable set of sentences of one \(\mathscr L(\tau)\) whose finite
      subsets all have models has a model.
\item \(\mathscr L(Q^{\mathrm H})\) is \emph{axiomatizable} if, for each
      finite vocabulary \(\tau\), the set of codes of valid sentences of
      \(\mathscr L(Q^{\mathrm H})(\tau)\) is recursively enumerable, for a
      coding of formulas by natural numbers under which the formation
      rules are computable and, for \(\tau\supseteq\Lang_A\), the passage
      from the code \(\ulcorner\theta\urcorner\) of a first-order
      \(\Lang_A\)-sentence (\cref{def:language-arithmetic}) to the code of
      \(\sigma\to\theta\), for a fixed sentence \(\sigma\), is computable.
\item For an ordinal \(\alpha\), \(Q_\alpha x\,\varphi(x)\) holds iff at
      least \(\aleph_\alpha\) elements satisfy \(\varphi\).  So \(Q_0\) is
      ``there exist infinitely many'' as in
      \cref{thm:p03-w03-p03-042-mundici-q-zero}, and \(Q_1\) is the
      quantifier of \cref{def:p03-q-one}.
\end{enumerate}
\end{hbconvention}

Walkoe observed that \(Q^{\mathrm H}\) is the weakest partially ordered
quantifier: if \(\mathscr L(Q)\not\equiv\mathscr L_{\omega\omega}\) for a
partially ordered quantifier \(Q\), then \(\mathscr L(Q^{\mathrm
H})\leq\mathscr L(Q)\) \cite[p.~213]{Mundici1985OtherQuantifiers}.  For
the quantifiers \(Q_\pi\) of \cref{def:p03-w03-p03-042-mundici-po-quantifier}
the hypothesis is a condition on the number of rows.

\begin{hbtheorem}{thm:p03-w03-p03-042-mundici-walkoe}{proved=annex:p03-w03-p03-042-mundici-walkoe}{Walkoe's observation for the quantifiers \(Q_\pi\)}
Let \(\pi\) be a pattern with \(r\) rows.
\begin{enumerate}
\item If \(r=1\), then \(\mathscr L(Q_\pi)\equiv\mathscr L_{\omega\omega}\).
\item If \(r\geq2\), then \(\mathscr L(Q^{\mathrm H})\leq\mathscr
      L(Q_\pi)\) and \(\mathscr L(Q_\pi)\not\equiv\mathscr L_{\omega\omega}\).
\end{enumerate}
Hence \(\mathscr L(Q_\pi)\not\equiv\mathscr L_{\omega\omega}\) iff
\(r\geq2\), and in that case \(\mathscr L(Q^{\mathrm H})\leq\mathscr
L(Q_\pi)\).  Part~1 uses the axiom of choice.
\end{hbtheorem}

The two theorems of Mundici's \S1.2 measure how far \(Q^{\mathrm H}\)
alone can replace the countably many partially ordered quantifiers, and how
partially ordered quantifiers relate to second-order logic; Mundici notes
that by a result of V\"a\"an\"anen no generalized quantifier \(Q\) has
\(\mathscr L(Q)\) equivalent to full second-order logic
\cite[p.~214]{Mundici1985OtherQuantifiers}.  The first theorem is proved
for an arbitrary theory with a definable pairing function, which is
Mundici's remark that the proof needs nothing else
\cite[p.~214]{Mundici1985OtherQuantifiers}.

\begin{hbtheorem}{thm:p03-w03-p03-042-mundici-pairing}{proved=annex:p03-w03-p03-042-mundici-pairing}{Contracting partially ordered prefixes with a pairing function (Mundici, Remark after Theorem~1.2.1)}
Let \(T\) be a first-order theory of vocabulary \(\tau_0\) and
\(\rho(x,y,z)\) a first-order \(\tau_0\)-formula that in every model
\(\mathfrak M\) of \(T\) defines the graph of a bijection from \(M^2\)
onto \(M\).  Let \(\tau\supseteq\tau_0\).  For every formula \(\varphi\) of
\(\mathscr L^{\circledR}(\tau)\) there is a formula \(\psi\) of
\(\mathscr L(Q^{\mathrm H})(\tau)\) with the same free variables such that
\(\varphi\leftrightarrow\psi\) holds, under every assignment, in every
\(\tau\)-structure whose \(\tau_0\)-reduct is a model of \(T\).  The axiom
of choice is used.
\end{hbtheorem}

\begin{hbcorollary}{thm:p03-w03-p03-042-mundici-peano}{proved=annex:p03-w03-p03-042-mundici-peano}{\(\mathscr L^{\circledR}\) and \(\mathscr L(Q^{\mathrm H})\) over Peano arithmetic (Mundici, Theorem~1.2.1)}
\(\mathscr L^{\circledR}\) is equivalent to \(\mathscr L(Q^{\mathrm H})\)
in first-order Peano arithmetic: for every \(\varphi\) in
\(\mathscr L^{\circledR}\) there is a \(\psi\) in \(\mathscr L(Q^{\mathrm
H})\) having the same models as \(\varphi\) among the models of Peano
arithmetic \cite[p.~214]{Mundici1985OtherQuantifiers}.  Here the
vocabulary of \(\varphi\) contains \(\Lang_A\), and a model of Peano
arithmetic is a structure whose \(\Lang_A\)-reduct satisfies the axioms
\(\PA\) of \cref{tab:p05-q-pa-axioms}.
\end{hbcorollary}

Mundici's proof first contracts each row to one universally and one
existentially quantified variable with the pairing function, then trades
\(n\) such rows for two rows of \(n\) variables each by the equivalence
\[
 \begin{Bmatrix}\forall x_1\ \exists z_1\\ \cdots\quad\cdots\\
 \forall x_n\ \exists z_n\end{Bmatrix}\theta
 \quad\text{and}\quad
 \begin{Bmatrix}\forall x_1\cdots x_n\ \exists z_1\cdots z_n\\
 \forall x'_1\cdots x'_n\ \exists z'_1\cdots z'_n\end{Bmatrix}
 \Bigl[\theta\wedge\bigwedge_{i=1}^n(x_i=x'_i\to z_i=z'_i)\Bigr],
\]
and contracts the two rows into \(Q^{\mathrm H}\)
\cite[p.~214]{Mundici1985OtherQuantifiers}.  The second row of the
right-hand prefix only checks that the \(i\)th value of the first row's
function depends on \(x_i\) alone.

\begin{hbexample}{def:p03-w03-p03-042-mundici-pairing-example}{Three rows over Peano arithmetic}
Take
\[
 \varphi=\begin{Bmatrix}\forall x_1\ \exists z_1\\ \forall x_2\ \exists
 z_2\\ \forall x_3\ \exists z_3\end{Bmatrix}\theta .
\]
In a model of \(\PA\) it says that there are
\(f_1,f_2,f_3:M\to M\) such that, for all \(x_1,x_2,x_3\),
\[
 \theta(x_1,x_2,x_3,f_1(x_1),f_2(x_2),f_3(x_3)).
\]  The two-row form asks for
\(F,F':M^3\to M^3\) with \(\theta(\bar x,F(\bar x))\) and \(F(\bar
x)_i=F'(\bar x')_i\) whenever \(x_i=x'_i\); taking \(\bar x'=\bar x\) and
then \(\bar x'\) arbitrary with \(x'_i=x_i\) shows that \(F(\bar x)_i\)
depends on \(x_i\) only, so \(f_i(a)=F(a,a,a)_i\) recovers the first
form.  Coding \(\bar x\) and \(F(\bar x)\) by the pairing graph
\(\mathsf P\) of \cref{lem:p05-pa-pair-beta} as single elements \(u\) and
\(v\) gives a formula \(Q^{\mathrm H}u\,u'\,v\,v'\,\theta^\ast\) of
\(\mathscr L(Q^{\mathrm H})\), with \(\theta^\ast\) first-order in
\(\theta\) and \(\mathsf P\).
\end{hbexample}

Mundici's \S1.2 continues with \(\Delta\mathscr L(Q^{\mathrm H})=\Delta
\mathscr L^{\circledR}=\Delta\mathscr L^{\mathrm{mII}}\) (his
Theorem~1.2.2), proved by two relative-projective-class reductions, with
the consequences on L\"owenheim and Hanf numbers and on the sets of valid
sentences recorded in his Corollary~1.2.3 and Remark~1.2.4
\cite[pp.~214--216]{Mundici1985OtherQuantifiers}.  These rest on the
notions \(\leq_{\mathrm{RPC}}\) and \(\Delta\)-closure of the general
framework, which this chapter does not set up; no statement of them is
made here.

Turning to the explicit expressive power of \(\mathscr L(Q^{\mathrm H})\)
\cite[p.~216]{Mundici1985OtherQuantifiers}: since \(Q_0\) is definable in
it (\cref{thm:p03-w03-p03-042-mundici-q-zero}), one sentence pins down the
standard model of arithmetic.

\begin{hbtheorem}{thm:p03-w03-p03-042-mundici-not-compact}{proved=annex:p03-w03-p03-042-mundici-not-compact}{\(\mathscr L(Q^{\mathrm H})\) is neither \((\omega,\omega)\)-compact nor axiomatizable (Mundici, Proposition~1.3.1)}
\begin{enumerate}
\item Let \(\sigma_{\mathbb N}\) be the conjunction of the seven axioms of
      \(\mathsf Q\) (\cref{tab:p05-q-pa-axioms}), of
      \(\forall x\forall y\,(x+y=y+x)\), and of \(\forall x\,\neg Q_0y\,
      \exists z\,(y+z=x)\), with \(Q_0\) replaced by its definition in
      \cref{thm:p03-w03-p03-042-mundici-q-zero}.  Then
      \(\sigma_{\mathbb N}\in\mathscr L(Q^{\mathrm H})(\Lang_A)\), and an
      \(\Lang_A\)-structure is a model of \(\sigma_{\mathbb N}\) iff it is
      isomorphic to \(\mathcal N\).
\item \(\mathscr L(Q^{\mathrm H})\) is not \((\omega,\omega)\)-compact.
\item \(\mathscr L(Q^{\mathrm H})\) is not axiomatizable.
\end{enumerate}
The axiom of choice enters through
\cref{thm:p03-w03-p03-042-mundici-q-zero}.
\end{hbtheorem}

Mundici's proof of Proposition~1.3.1 cites G\"odel's incompleteness
theorem for the second item; the annex uses Tarski's undefinability
theorem (\cref{thm:tarski-undefinability}) as in the proof of
\cref{thm:p03-w03-p03-042-ring-lemma}.  The same proposition also states
that \(\mathscr L(Q^{\mathrm H})\) lacks the weak Beth property, which
Mundici obtains as a case of a theorem of abstract model theory for
finitely generated logics in which the class of structures isomorphic to
\(\langle\omega,<\rangle\) is EC, citing Makowsky and Shelah
\cite[p.~216]{Mundici1985OtherQuantifiers}; that part is not proved in
this chapter.

\begin{hbtheorem}{thm:p03-w03-p03-042-mundici-pure-identity}{proved=annex:p03-w03-p03-042-mundici-pure-identity}{Pure identity formulas of \(\mathscr L(Q^{\mathrm H})\) on infinite sets}
For every formula \(\varphi\) of \(\mathscr L(Q^{\mathrm H})\) of the
empty vocabulary (only \(=\) occurs) there is a first-order formula
\(\hat\varphi\) of the empty vocabulary with the same free variables such
that \(\kappa\models\varphi\leftrightarrow\hat\varphi\), under every
assignment, for every set \(\kappa\) with \(\kappa\geq\omega\).  The axiom
of choice is used.
\end{hbtheorem}

\begin{hbtheorem}{thm:p03-w03-p03-042-mundici-q-alpha}{proved=annex:p03-w03-p03-042-mundici-q-alpha}{Which \(Q_\alpha\) the Henkin quantifier defines (Mundici, Theorem~1.3.2)}
\(\mathscr L(Q^{\mathrm H})\geq\mathscr L(Q_\alpha)\) iff \(\alpha=0\)
\cite[p.~216]{Mundici1985OtherQuantifiers}.
\end{hbtheorem}

Mundici reduces the direction from left to right to
\cref{thm:p03-w03-p03-042-mundici-pure-identity}, whose only nontrivial
case is a formula \(Q^{\mathrm H}\psi\), handled with the upward and
downward L\"owenheim--Skolem theorems for \(\mathscr L_{\omega\omega}\);
for \(\alpha>0\), \(Q_\alpha\) separates infinite sets
\cite[p.~216]{Mundici1985OtherQuantifiers}.  In particular \(Q_1\) of
\cref{def:p03-q-one} is not definable in \(\mathscr L(Q^{\mathrm H})\),
although \(\mathscr L(Q^{\mathrm H})\) defines \(Q_0\).

\begin{hbwarning}{def:p03-w03-p03-042-mundici-limits}{What the Henkin quantifier does not give}
\(\mathscr L(Q^{\mathrm H})\) gives up \((\omega,\omega)\)-compactness,
one of the two properties that Lindstr\"om's theorem
(\cref{thm:p03-lindstrom}) ties to first-order logic, and it has no
recursively enumerable set of valid sentences
(\cref{thm:p03-w03-p03-042-mundici-not-compact}).  Its gain in expressive
power is uneven: it defines \(Q_0\), and among the models of \(\PA\) it
defines every first-order projective class
(\cref{thm:p03-w03-p03-042-mundici-pc-ec,thm:p03-w03-p03-042-mundici-peano}),
yet on pure sets it sees no difference between infinite cardinalities
(\cref{thm:p03-w03-p03-042-mundici-pure-identity}).
\end{hbwarning}

\paragraph{Attributions recorded by Mundici.}  According to Mundici's
bibliographical notes \cite[p.~217]{Mundici1985OtherQuantifiers}:
Henkin introduced \(Q^{\mathrm H}\) in 1961, and Ehrenfeucht proved that
\(\mathscr L(Q^{\mathrm H})\) is neither countably compact nor
axiomatizable; Theorems~1.1.2 and~1.2.1 are proved in Enderton (1970) and
Walkoe (1970); the proof of Theorem~1.3.2 is due to Lopez-Escobar (1969),
who also proved that interpolation fails; Paulos (1976) proved that the
\(\Delta\)-closure and Beth properties fail for \(\mathscr L(Q^{\mathrm
H})\), and Gostanian and Hrbacek (1976) proved the failure of the weak
Beth property; Theorem~1.2.2 is proved in Krynicki (1978) and Krynicki and
Lachlan (1979), the latter also containing decidability and
undecidability results on \(\mathscr L(Q^{\mathrm H})\); Krynicki (1977b)
uses back-and-forth games for \(\mathscr L(Q^{\mathrm H})\)-equivalence
in connection with Theorem~1.3.2, with related work by Krawczyk and
Krynicki (1976) and Weese (1980); Kreisel (1967), Mostowski (1968) and
Lindstr\"om (1969) are cited for more on these questions.  Mundici further records two uses of
partially ordered quantifiers: Barwise (1976) obtains first-order
axiomatizations of classes such as the structures with a nontrivial
automorphism \(f\) with \(f^2\) the identity; and Walkoe (1970, 1976) and
Keisler and Walkoe (1973) prove that for distinct first-order prefixes
\(Q'\) and \(Q''\) of the same length there is a quantifier-free
\(\varphi\) of \(\mathscr L_{\omega\omega}\) for which no quantifier-free
\(\psi\) of \(\mathscr L_{\omega\omega}\) makes \(Q'\varphi\) equivalent to
\(Q''\psi\).  Harel (1979), Cowles (1981) and Barwise (1979) are given as
further reading on \(Q^{\mathrm H}\).  These results are reported here,
not proved.


\section{The effective content of Keisler's theorem for \(L(Q)\)}
\label{sec:p03-w03-p03-042-keisler-effective}

Keisler's logic \(L(Q)\), its weak and standard models and its axioms and
rules are \cref{def:p02-w02c-p02-003-kq-language,def:p02-w02c-p02-003-kq-calculus};
his completeness theorem, with the cardinality bounds of his
Corollary~2.9 and the first half of his Corollary~2.10, is
\cref{thm:p02-w02c-p02-003-keisler-completeness}, and his compactness
theorem is \cref{thm:p02-w02c-p02-003-keisler-compactness}.  The second
half of Corollary~2.10 states that the set of all valid sentences of
\(L(Q)\) is recursively enumerable in the set of symbols of \(L\)
\cite[\S2, p.~26]{Keisler1970}; Keisler gives no separate argument, and
records that Vaught first proved it by other methods
\cite[\S2, p.~27]{Keisler1970}.  The hypotheses under which ``recursive
in the set of symbols'' has a definite meaning are fixed first.

\begin{hbdefinition}{def:p03-w03-p03-042-presentation}{A presentation of a countable language}
Let \(L\) be a first-order language with identity and at most countably
many predicate, function and individual constant symbols.  A
\emph{presentation} of \(L\) is an injective assignment of natural numbers
to the symbols of \(L\), together with the set \(\Sigma_L\subseteq\mathbb
N\) of numbers so assigned and the function \(\kappa_L\) giving, for each
\(e\in\Sigma_L\), the kind (predicate, function or constant) and the
number of places of the symbol numbered \(e\).  The logical symbols
\(\neg,\land,\exists,\forall,Q,=\), the parentheses and the variables
\(v_0,v_1,\ldots\) receive fixed numbers outside the range of the
assignment.  Expressions are coded by the numbers of finite sequences of
symbol numbers.  A set of expressions is \emph{recursive}
(\emph{recursively enumerable}) \emph{in the set of symbols of \(L\)} if
its set of codes is recursive (recursively enumerable) relative to the
oracle \((\Sigma_L,\kappa_L)\).
\end{hbdefinition}

\begin{hbtheorem}{thm:p03-w03-p03-042-kq-effective}{proved=annex:p03-w03-p03-042-kq-effective}{The axioms of \(L(Q)\) are recursive and the valid sentences are recursively enumerable}
Let \(L\) be as in \cref{def:p03-w03-p03-042-presentation}, with a fixed
presentation.
\begin{enumerate}
\item The set of formulas of \(L(Q)\), the set of sentences, and the set
      of axioms of \cref{def:p02-w02c-p02-003-kq-calculus} (Axiom~0 in the
      form (0a)--(0g), and Axioms~1--4) are recursive in the set of symbols
      of \(L\).
\item The relation ``\(p\) is a deduction of \(\varphi\) from the empty
      set'' (a finite sequence of formulas, each an axiom or obtained from
      earlier ones by modus ponens or generalization, ending in
      \(\varphi\)) is recursive in the set of symbols of \(L\).
\item The set of provable formulas of \(L(Q)\), and the set of valid
      sentences of \(L(Q)\), are recursively enumerable in the set of
      symbols of \(L\) \cite[\S2, Corollary~2.10, p.~26]{Keisler1970}.
\end{enumerate}
In particular, if \(L\) is finite, or if \(\Sigma_L\) and \(\kappa_L\) are
recursive, the set of valid sentences of \(L(Q)\) is recursively
enumerable.
\end{hbtheorem}

The proof in the annex describes decision procedures and identifies them
with procedures recursive in the oracle through the Church--Turing thesis
of \cref{def:p04-church-turing}; that identification is the only
non-mathematical step.  Validity here is truth in all \emph{standard}
models \cite[\S2, p.~26]{Keisler1970}; the set of sentences true in all
weak models in which the axioms hold is the same set, by
\cref{thm:p02-w02c-p02-003-kq-weak-completeness}.

\begin{hbwarning}{def:p03-w03-p03-042-axiom-zero}{Effectivity and the choice of Axiom~0}
Keisler's Axiom~0 consists of the first-order axiom schemes of Church's
textbook, which he cites but does not list
\cite[\S1, p.~6]{Keisler1970}; the handbook's calculus uses (0a)--(0g)
instead (\cref{def:p02-w02c-p02-003-kq-axiom-zero}).  The effectivity
theorem above is proved for (0a)--(0g).  The set of valid sentences is
defined semantically and does not depend on any list of axioms, so the
recursive enumerability of the valid sentences holds as stated.  Items~1
and~2, on the axiom set and on deductions, are established here for
(0a)--(0g) only; Church's list is not reproduced in the handbook and has
not been compared with (0a)--(0g).
\end{hbwarning}

\section{From Mostowski's and Keisler's notation to the handbook's}
\label{sec:p03-w03-p03-042-bridge}

\begin{hbconvention}{def:p03-w03-p03-042-mostowski-bridge}{Mostowski's quantifiers as isomorphism-invariant classes}
Let \(Q\) be an unlimited quantifier in Mostowski's sense.  Its
\emph{class form} is
\[
 \mathsf K(Q)=\{(M,A):M\neq\varnothing,\ A\subseteq M,\ Q_M(A)=\vee\},
\]
where \(A\) is identified with a one-argument propositional function on
\(M\).  Conversely a unary generalized quantifier \(\mathcal K\) of
\cref{def:p02-generalized-quantifier} gives the unlimited quantifier
\(Q^{\mathcal K}\) with \(Q^{\mathcal K}_M(A)=\vee\) iff
\((M,A)\in\mathcal K\), for nonempty \(M\).
\end{hbconvention}

\begin{hbtheorem}{thm:p03-w03-p03-042-bridge}{proved=annex:p03-w03-p03-042-bridge}{Correspondence between the two presentations}
\begin{enumerate}
\item \(\mathsf K(Q)\) is an isomorphism-closed class of pairs, hence a
      unary generalized quantifier in the sense of
      \cref{def:p02-generalized-quantifier}, and \(Q\mapsto\mathsf K(Q)\)
      is a bijection between Mostowski's unlimited quantifiers restricted
      to nonempty sets and unary generalized quantifiers, with inverse
      \(\mathcal K\mapsto Q^{\mathcal K}\).
\item For every formula \(\mathrm Z\) of (S) and every nonempty \(I\), let
      \(\mathcal M\) be the structure with universe \(I\) interpreting a
      \(k\)-ary relation symbol \(R_j\) by
      \(\{\bar a\in I^k:[\mathrm F_j]_M(\bar a,a_k,a_k,\ldots)=\vee\}\) and
      let \(s(\mathrm x_i)=[\mathrm x_i]_M\).  Then
      \(\mathit{val}_{MI}(\mathrm Z)=\vee\) iff
      \(\mathcal M,s\entails\mathrm Z^\circ\) in
      \(\mathsf{FOL}(\{\mathsf K(Q^1),\ldots,\mathsf K(Q^s)\})\), where
      \(\mathrm Z^\circ\) replaces \(\mathrm F_j\) by \(R_j\), the stroke
      by \(\neg(\cdot\land\cdot)\), and \(\mathbf Q^j\) by the quantifier
      symbol for \(\mathsf K(Q^j)\).
\item The uncountability quantifier of \cref{def:p03-q-one} is
      \(Q_1=\mathsf K({\sim}P)\): it is determined, on each \(I\), by
      \(T(\mathfrak m,\mathfrak n)=\vee\) iff \(\mathfrak m>\aleph_0\).
      It is not numerical on any infinite \(I\) with
      \(\overline I>\aleph_0\), it has value \(\wedge\) everywhere on a
      denumerable \(I\), and it does not have property (E).
\end{enumerate}
\end{hbtheorem}

The two presentations differ in three respects.  First,
\cref{def:p02-generalized-quantifier} allows \(n\)-ary quantifiers binding
a tuple of variables, while Mostowski's bind one variable; part~1 is the
unary case.  Second, Mostowski requires invariance under one-one maps
between different sets as part of the notion of an unlimited quantifier,
and invariance under permutations for a limited one; the handbook builds
both into isomorphism closure.  Third, (S) has free functional variables
and defines truth as truth under all valuations of them, so ``true''
formulas of (S) are the second-order universal closures of
\(\mathsf{FOL}\) formulas, and the completeness problem for (S) is the
problem of enumerating the valid sentences of
\(\mathsf{FOL}(\{\mathsf K(Q^1),\ldots,\mathsf K(Q^s)\})\) in a countable
relational signature (part~2).  The handbook requires \(M\neq\varnothing\); Mostowski
places no such restriction on \(I\), so the correspondence concerns
nonempty sets.

\begin{hbconvention}{def:p03-w03-p03-042-keisler-bridge}{Keisler's \(L(Q)\) in the handbook's notation}
Keisler's \(Q\), read in standard models, is the handbook's \(Q_1\):
\((\mathfrak A,q)\) with \(q\) the set of uncountable subsets of \(A\) is
the structure \(\mathfrak A\) with the satisfaction clause of
\cref{def:p03-q-one}, so \(L(Q)\) with standard models is
\(\mathsf{FOL}(\{Q_1\})\) over the signature \(L\), and
\(\mathfrak A\models\varphi[a_1,\ldots,a_n]\) is
\(\mathfrak A,s\entails\varphi\) with \(s(v_i)=a_i\).  Keisler writes
identity as a triple bar; it is written \(=\).  Keisler's ``countable''
means of power at most \(\omega\), and his \(\omega_1\) is the cardinal
\(\aleph_1\).  Weak models have no counterpart among the handbook's
structures: they are an auxiliary semantics in which \(Q\) is interpreted
by an arbitrary family \(q\) of subsets; they serve only in the proof.
\end{hbconvention}

\begin{hbtheorem}{thm:p03-w03-p03-042-q-one-host-form}{proved=annex:p03-w03-p03-042-q-one-host-form}{The host statement on \(L(Q_1)\), proved}
Let \(\tau\) be a signature and \(L(Q_1)\) the logic
\(\mathsf{FOL}(\{Q_1\})\) over \(\tau\).
\begin{enumerate}
\item For every sentence \(\sigma\) of \(L(Q_1)\): \(\sigma\) is valid iff it
      is derivable in Keisler's calculus
      (\cref{def:p02-w02c-p02-003-kq-calculus}) over \(\tau\), whatever the
      cardinality of \(\tau\).
\item Every countable set \(\Sigma\) of \(L(Q_1)\)-sentences whose finite
      subsets are satisfiable has a model, whatever the cardinality of
      \(\tau\).
\item If \(\tau\) is countable, every satisfiable set of
      \(L(Q_1)\)-sentences has a model of cardinality at most \(\aleph_1\),
      and it has one of cardinality exactly \(\aleph_1\) iff
      \(\Sigma\cup\{Q_1x\,(x=x)\}\) is satisfiable.
\item If \(\tau\) contains a unary relation symbol and \(\aleph_1\)
      constant symbols, the conclusion of item~2 fails for some set
      \(\Sigma\) of power \(\aleph_1\), and so does completeness for sets of
      premises: there is a consistent \(\Sigma\) with no model.
\end{enumerate}
\end{hbtheorem}

Item~1 and item~2 are the two halves of \cref{thm:p03-keisler-q-one}:
its ``sound complete axiomatization'' is item~1, and its ``countable set
of sentences'' is a hypothesis on the set, not on the signature, which
item~2 shows to be enough.  Its ``countable signature'' bound is item~3,
whose second clause is the source's Corollary~2.9 in model-theoretic
form; the source's statements concern a countable \emph{language}
\cite[\S2, Corollary~2.9, p.~26]{Keisler1970}.  Item~4 is Keisler's
example \cref{def:p02-w02c-p02-003-kq-example}, which shows that the
countability hypothesis on \(\Sigma\) in item~2 cannot be dropped.

\begin{hbexample}{def:p03-w03-p03-042-worked}{A worked reconstruction: \(Q_1\) and \(P\) on three sets}
Take \(I_1=\{0,1,2\}\), \(I_2=\mathbb N\) and \(I_3=\omega_1\), and
\(A=\{x\in I:x\text{ is even}\}\) on the first two, \(A=\{\alpha<\omega_1:
\alpha\text{ is a limit ordinal or }0\}\) on the third.  The cardinal pairs
\((\overline A,\overline{I-A})\) are \((2,1)\), \((\aleph_0,\aleph_0)\) and
\((\aleph_1,\aleph_1)\).  Mostowski's \(P\) gives \(\vee,\vee,\wedge\)
(at most denumerably many), so \({\sim}P\) gives \(\wedge,\wedge,\vee\),
and by part~3 of \cref{thm:p03-w03-p03-042-bridge} these are the values
of \(Q_1x\,A(x)\).  The Theorem~1 description of \(Q_1\) on \(I_3\) needs
only the pairs with \(\mathfrak m+\mathfrak n=\aleph_1\); exactly those
with \(\mathfrak m=\aleph_1\) receive \(\vee\).  On \(I_2\), the function
\(T\) of \(Q_1\) is constantly \(\wedge\), so \(Q_1\) limited to a
denumerable set is the numerical quantifier \(\overset{(0)}{\Sigma}\wedge
{\sim}\overset{(0)}{\Sigma}\); Theorem~4 then gives a positive solution of
the completeness problem for \(\exists,\forall,Q_1\) limited to a
denumerable set.  This is consistent with Keisler's theorem: there the
models are not restricted to one cardinality, and \(Q_1\) is not numerical
on uncountable sets.
\end{hbexample}

\begin{hbconvention}{def:p03-w03-p03-042-monadic-bridge}{The monadic calculus in the handbook's notation}
Under \cref{thm:p03-w03-p03-042-bridge}(2), a formula of \((\mathrm
S^M_n)\) is an equality-free formula of
\(\mathsf{FOL}(\{\mathsf K(Q^1),\ldots,\mathsf K(Q^s)\})\) over the
signature of unary relation symbols \(R_1,\ldots,R_n\), with \(\mathrm
F_h\) read as \(R_h\); \((\mathrm S^M)\) uses \(R_1,R_2,\ldots\).
Mostowski's \(\mathrm F_h^0,\mathrm F_h^1\) are \(R_h,\neg R_h\); a
constituent of order \(n\) is a complete conjunction of \(R_h(x)\) and
\(\neg R_h(x)\) for \(h\leq n\), whose extensions are the \(2^n\) cells of
the structure.  The quantifier ``at least \(\mathfrak m\)'' of
Theorem~10 is the class \(\mathsf K_{\geq\mathfrak m}=\{(M,A):
M\neq\varnothing,\ A\subseteq M,\ \overline A\geq\mathfrak m\}\), so that
\(\mathsf K_{\geq1}\) is \(\exists\), \(\mathsf K_{\geq\aleph_0}=\mathsf
K({\sim}S)\) is \(Q_0\) (``there exist infinitely many''), and
\(\mathsf K_{\geq\aleph_1}=\mathsf K({\sim}P)\) is \(Q_1\) of
\cref{def:p03-q-one}.  Mostowski's ``true'' formula is a formula valid in
all structures under all assignments, and ``recursive'' is
``decidable''.
\end{hbconvention}

\begin{hbtheorem}{thm:p03-w03-p03-042-monadic-bridge}{proved=annex:p03-w03-p03-042-monadic-bridge}{Theorems~10 and~11 in the handbook's notation}
Let \(\tau\) be a signature of unary relation symbols \(R_1,R_2,\ldots\)
and let all formulas be equality-free.
\begin{enumerate}
\item For infinite cardinals \(\mathfrak m_2<\cdots<\mathfrak m_s\), the
      set of formulas of \(\mathsf{FOL}(\{\mathsf K_{\geq\mathfrak
      m_2},\ldots,\mathsf K_{\geq\mathfrak m_s}\})\) over \(\tau\) that are
      valid is decidable.  In particular the valid equality-free monadic
      formulas of \(\mathsf{FOL}(\{Q_0,Q_1\})\) form a decidable set.
\item For any unary generalized quantifiers \(\mathcal K^1,\ldots,\mathcal
      K^s\) and any \(n\geq1\), the set of valid formulas of
      \(\mathsf{FOL}(\{\mathcal K^1,\ldots,\mathcal K^s\})\) over
      \(\{R_1,\ldots,R_n\}\) is decidable.
\end{enumerate}
\end{hbtheorem}

Item~1 is uniform in the cardinals and comes with the procedure of
Lemmas~(c) and~(d); item~2 asserts only that a decision procedure exists,
as Mostowski's Fermat-prime example shows
(\cref{thm:p03-w03-p03-042-fermat}).  Mostowski's own route to
Theorem~10 is the elimination of nested quantifiers through additivity;
his route to Theorem~11 replaces additivity, which arbitrary quantifiers
lack, by quantification over every union of cells.  Both are kept in the
annex.

\begin{hbexample}{def:p03-w03-p03-042-monadic-worked}{A worked reconstruction for \S6}
\emph{Mycielski's lemma.}  Take \(A_1=\{1,2\}\) and \(A_2=\{1,2,8\}\),
subsets of \(A_0\) containing \(1\).  Then \(8=2^{3^1}\in A_2-A_1\), and
\(8\in H(A_2)\) because eight summands \(1\) give \(8\in A_2\); but
\(8\notin H(A_1)\), since eight summands from \(\{1,2\}\) add up to at
least \(8\), which is not in \(A_1\).

\emph{Theorem~9 with \(n=2\).}  The constituents of order \(2\) in
lexicographic order have indices \((0,0),(0,1),(1,0),(1,1)\), so
\(\mathrm S_1(\mathrm x)=\mathrm F_1(\mathrm x)\wedge\mathrm F_2(\mathrm
x)\) and \(\mathrm S_2(\mathrm x)=\mathrm F_1(\mathrm x)\wedge{\sim}
\mathrm F_2(\mathrm x)\).  In the handbook's notation \(\mathrm W_2\) says
that \(R_1\cap R_2\), \(R_1-R_2\) and \(R_1\) each have a power lying in
\(A\); it is satisfiable in \(\mathbb N\) iff two members of \(A\) have a
sum in \(A\), that is, iff \(2\in H(A)\).

\emph{Theorem~10 with \(s=2\), \(\mathfrak m_1=1\), \(\mathfrak
m_2=\aleph_1\).}  Here \(Q^1=\exists\) and \(Q^2={\sim}P\), which is
\(Q_1\).  Consider
\[
 \mathrm Z=(\mathbf Q^1\mathrm x)\mathrm F_1(\mathrm x)\supset
 [(\mathbf Q^2\mathrm x)\mathrm F_1(\mathrm x)\vee(\mathbf Q^2\mathrm x)
 {\sim}\mathrm F_1(\mathrm x)],
\]
in the handbook's notation \(\exists x\,R_1(x)\to(Q_1x\,R_1(x)\lor Q_1x\,
\neg R_1(x))\).  With \(n=1\) the constituents are \(\mathrm F_1(\mathrm
x)\) (index \(0\)) and \({\sim}\mathrm F_1(\mathrm x)\) (index \(1\)), so
Lemma~(c) returns \(\mathrm U=\mathrm A^1_0\supset(\mathrm A^2_0\vee
\mathrm A^2_1)\) and \(\Phi_{\mathrm U}=a^1_0\supset(a^2_0\vee a^2_1)\).
The pair \(\gamma(a^1_0)=\vee\), \(\gamma(a^2_0)=\gamma(a^1_1)=
\gamma(a^2_1)=\wedge\) is allowable (\(a^2_e\leq a^1_e\)) and admissible
(some \(a^1_e\) is \(\vee\); there are no free variables), and gives
\(\Phi_{\mathrm U}\) the value \(\wedge\); so \(\mathrm Z\) is not true.
The construction in the proof of Lemma~(d) gives \(j_0=1\), \(j_1=0\),
hence a cell of one element and an empty cell: \(I=\{p\}\) with
\([\mathrm F_1]=\{p\}\), where indeed something is \(\mathrm F_1\) but
neither \(\mathrm F_1\) nor its complement is uncountable.  For
\(\mathrm Z'=(\mathbf Q^2\mathrm x)\mathrm F_1(\mathrm x)\supset(\mathbf
Q^1\mathrm x)\mathrm F_1(\mathrm x)\), the polynomial \(a^2_0\supset
a^1_0\) has the value \(\vee\) at every allowable valuation, since
\(\gamma(a^2_0)\leq\gamma(a^1_0)\); so \(\mathrm Z'\), that is
\(Q_1x\,R_1(x)\to\exists x\,R_1(x)\), is true.
\end{hbexample}
